% =====================================================================
% Paper 48 (standalone): Krein-pointed proper quantum metric spaces and
% the bridge to Mondino-Sämann Lorentzian pre-length spaces, with
% application to the GeoVac hemispheric wedge.
%
% Tenth math.OA standalone in the GeoVac series. Siblings:
%   Paper 29 (Ramanujan-Hopf), Paper 32 (spectral triple), Paper 38
%   (SU(2) Riemannian propinquity), Paper 39 (tensor-product propinquity),
%   Paper 40 (universal rate constant on compact Lie groups), Paper 42
%   (Tomita-Takesaki Riemannian closure), Paper 43 (Lorentzian extension
%   of the four-witness theorem), Paper 44 (Lorentzian operator-system
%   substrate), Paper 45 (K+-weak-form Lorentzian propinquity
%   convergence), Paper 46 (strong-form Lorentzian propinquity
%   convergence on the natural and enlarged substrates), Paper 47
%   (two-rate hybrid convergence to the non-compact carrier).
%
% Status: first draft, May 2026. Built from Phase A.5' synthesis
% (debug/sprint_phase_a5prime_synthesis_paper48_outline_memo.md) and the
% six Phase A sub-sprint memos (Tier 3-Light, A.2', A.3' + concurrent-
% work re-check, A.4'-A, A.4'-B, A.4'-C).
%
% Main result: the Wick-rotation functor W: KreinMetaMet_pp ->
% LorPLG_cov, identified via R3 (Connes-Rovelli thermal time) as the
% structural bridge resolving the F2 forward-vs-reverse triangle
% mismatch between Latremoliere 2512.03573 pinned proper QMS framework
% and Mondino-Samann arXiv:2504.10380 v4 covered Lorentzian pre-length
% space framework. Application to the GeoVac hemispheric wedge yields
% seven newly accessible theorems including T6 (G2 metric-equivalent
% closure for the wedge, a published-open-problem workaround at the
% wedge-specific level).
%
% Honest scope: K+-weak-form bridge only. Strong-form bridge with
% enlarged substrate (Q1') deferred as a multi-month NCG-mathematics
% follow-on.
% =====================================================================
\documentclass[11pt,a4paper]{article}

\usepackage[T1]{fontenc}
\usepackage[utf8]{inputenc}
\usepackage{amsmath, amssymb, amsthm, mathrsfs, mathtools}
\usepackage{newunicodechar}
\newunicodechar{ε}{\ensuremath{\varepsilon}}
\newunicodechar{α}{\ensuremath{\alpha}}
\newunicodechar{β}{\ensuremath{\beta}}
\newunicodechar{γ}{\ensuremath{\gamma}}
\newunicodechar{δ}{\ensuremath{\delta}}
\newunicodechar{κ}{\ensuremath{\kappa}}
\newunicodechar{λ}{\ensuremath{\lambda}}
\newunicodechar{μ}{\ensuremath{\mu}}
\newunicodechar{π}{\ensuremath{\pi}}
\newunicodechar{σ}{\ensuremath{\sigma}}
\newunicodechar{τ}{\ensuremath{\tau}}
\newunicodechar{ω}{\ensuremath{\omega}}
\newunicodechar{Δ}{\ensuremath{\Delta}}
\newunicodechar{Λ}{\ensuremath{\Lambda}}
\newunicodechar{Σ}{\ensuremath{\Sigma}}
\newunicodechar{Ω}{\ensuremath{\Omega}}
\newunicodechar{Đ}{\ensuremath{\text{\DJ}}}
\usepackage{geometry}
\geometry{margin=1in}
\usepackage{hyperref}
\hypersetup{colorlinks=true, linkcolor=blue!50!black, citecolor=blue!50!black, urlcolor=blue!50!black}
% \usepackage{microtype}  % disabled: MiKTeX font-expansion environmental issue
\usepackage{graphicx}
\usepackage{booktabs}
\usepackage[numbers,sort&compress]{natbib}
\usepackage{xcolor}

\theoremstyle{plain}
\newtheorem{theorem}{Theorem}[section]
\newtheorem{lemma}[theorem]{Lemma}
\newtheorem{proposition}[theorem]{Proposition}
\newtheorem{corollary}[theorem]{Corollary}
\theoremstyle{definition}
\newtheorem{definition}[theorem]{Definition}
\newtheorem{remark}[theorem]{Remark}
\newtheorem{example}[theorem]{Example}

\newcommand{\nmax}{n_{\max}}
\newcommand{\Nt}{N_{t}}
\newcommand{\sthree}{S^{3}}
\newcommand{\Sone}{S^{1}}
\newcommand{\SoneT}{S^{1}_{T}}
\newcommand{\Manifold}{S^{3}\times S^{1}_{T}}
\newcommand{\Hilb}{\mathcal{H}}
\newcommand{\Krein}{\mathcal{K}}
\newcommand{\Op}{\mathcal{O}}
\newcommand{\Acal}{\mathcal{A}}
\newcommand{\Tcal}{\mathcal{T}}
\newcommand{\Bcal}{\mathcal{B}}
\newcommand{\Scal}{\mathcal{S}}
\newcommand{\Mcal}{\mathcal{M}}
\newcommand{\Ucal}{\mathcal{U}}
\newcommand{\Dcal}{\mathcal{D}}
\newcommand{\opnorm}[1]{\lVert #1 \rVert_{\mathrm{op}}}
\newcommand{\Frob}[1]{\lVert #1 \rVert_{F}}
\newcommand{\cbnorm}[1]{\lVert #1 \rVert_{\mathrm{cb}}}
\newcommand{\norm}[1]{\lVert #1 \rVert}
\newcommand{\R}{\mathbb{R}}
\newcommand{\C}{\mathbb{C}}
\newcommand{\N}{\mathbb{N}}
\newcommand{\Z}{\mathbb{Z}}
\newcommand{\Q}{\mathbb{Q}}
\newcommand{\SU}{\mathrm{SU}}
\newcommand{\Uone}{U(1)}
\newcommand{\Tr}{\mathrm{Tr}}
\newcommand{\DGV}{D_{\mathrm{GV}}}
\newcommand{\DL}{D_{L}}
\newcommand{\JL}{J_{L}}
\newcommand{\HGV}{\mathcal{H}_{\mathrm{GV}}}
\newcommand{\Kplus}{\Krein^{+}}
\newcommand{\Kminus}{\Krein^{-}}
\newcommand{\spat}{\mathrm{spat}}
\newcommand{\temp}{\mathrm{temp}}
\newcommand{\Bjoint}{B^{\mathrm{joint}}}
\newcommand{\Pjoint}{P^{\mathrm{joint}}}
\newcommand{\gammajoint}{\gamma^{\mathrm{joint}}}
\newcommand{\Cthreejoint}{C_{3}^{\mathrm{joint}}}
\newcommand{\Lprop}{\Lambda^{L}}
\newcommand{\Dethmet}{\eth^{K}}
\newcommand{\BigDeth}{\text{\DJ}^{K}}
\newcommand{\Wfun}{\mathbf{W}}
\newcommand{\KreinMM}{\mathbf{KreinMetaMet}_{\mathrm{pp}}}
\newcommand{\LorPLG}{\mathbf{LorPLG}_{\mathrm{cov}}}
\newcommand{\BWvac}{\omega_{W}^{L}}
\newcommand{\Kalpha}{K_{\alpha}^{W}}
\newcommand{\Hloc}{H_{\mathrm{local}}}
\newcommand{\Dwedge}{D_{W}^{L}}
\newcommand{\tmod}{\tau_{\mathrm{mod}}}
\newcommand{\kappag}{\kappa_{g}}

\title{Krein-pointed proper quantum metric spaces and the bridge to\\
Mondino--S\"amann Lorentzian pre-length spaces, with application to the\\
GeoVac hemispheric wedge}
\author{J.~Loutey\thanks{Independent researcher. \texttt{jloutey@gmail.com}.
This work was developed in an AI-augmented agentic workflow with
Anthropic's Claude. Implementation, internal memos, and bibliographic
collation were performed collaboratively; all theorems, design choices,
and editorial decisions are the author's responsibility.}}
\date{May 2026 (preprint)}

\begin{document}
\maketitle

\begin{abstract}
We construct a categorical bridge between Latr\'emoli\`ere's pinned
proper quantum metric space (PPQMS) framework~\cite{latremoliere2025_hypertopology}
and Mondino--S\"amann's covered Lorentzian pre-length space (LPLS)
framework~\cite{mondino_samann2025_pointed} (descoped; see the Status note: the bridge's domain category is
instantiated on the Paper~45 substrate, whose Lipschitz seminorm is
degenerate, so the metric-level theorems are open pending repair
and the bridge survives as categorical design). The bridge is a contravariant
functor
\[
\Wfun: \KreinMM \longrightarrow \LorPLG
\]
sending a Krein-pointed proper QMS $(\Acal^{K}, L^{K}, \Mcal^{L}, \BWvac)$
to a covered LPLS $(\hat{\Mcal}^{L}, \ell^{L}, \hat{\omega}_{W}^{L},
\hat{\Ucal}^{L})$ via Gelfand spectrum of the M-diagonal Abelian
topography, Wick-rotated modular flow as time separation, the
Bisognano--Wichmann vacuum as basepoint, and admissible-scaling
truncated covers. The structural mismatch between Latr\'emoli\`ere's
4-point relaxed forward metametric and Mondino--S\"amann's 3-point
reverse triangle (the F2 mismatch) is resolved via the Connes--Rovelli
thermal-time hypothesis~\cite{connes_rovelli1994}:\ the Krein
metametric measures \emph{thermal time} of the wedge KMS state, the
synthetic time separation measures \emph{geometric proper time}, and
the Wick rotation between them effects the analytic continuation that
maps the additive (forward-triangle) thermal-time structure to the
super-additive (reverse-triangle) geometric-time structure.

\medskip

\noindent\emph{Main theorem.} The Wick-rotation functor $\Wfun$ is
well-defined and the Krein--Mondino--S\"amann Bridge Theorem
(Theorem~\ref{thm:bridge}) holds at theorem-grade rigor for the
\emph{structural} legs --- (B1) structural correspondence on $15$ rows
and (B3) pre-compactness inheritance via Paper~44~\cite{paper44}
propagation-number bound --- while the metric-level legs ride Paper~45's
degenerate Krein seminorm:\ (B2) the reverse triangle inequality holds
on-orbit via modular-flow equality and is \emph{structurally vacuous
off-orbit} (the M-diagonal topography forces orbit-pair contradiction),
a structural/categorical property rather than a metric measurement; and
(B4) convergence transport via Plancherel-nested Berezin compatibility
with the Mondino--S\"amann Def~3.6 axioms is \textbf{descoped} ---
conditional pending substrate repair, not an established metric
convergence (see the Status note).

\medskip

\noindent\emph{GeoVac wedge application.} The bridge applies cleanly
to the Paper~43~\cite{paper43} hemispheric wedge at every finite
cutoff. Seven newly accessible theorems follow
(Theorems~\ref{thm:wedge_as_lpls}--\ref{thm:g2_wedge_closure} +
Corollary~\ref{cor:three_carrier_synthetic}), headlined by
\textbf{T6} (\emph{descoped}; see the Status note):\ a \emph{proposed}
$\mathbf{G2}$ metric-level closure for the GeoVac wedge,
\emph{conditional on the upstream metric repair} --- it inherits
Paper~45's Krein-side $\Lprop$, which is a degenerate seminorm (an
evaluation of a closed-form rate formula, not a distance), so T6 is not
yet a metric-level closure. The synthetic rate panel transports the
bit-exact values $\{2.0746, 1.6101, 1.3223\}$ from that formula
(rate-formula values, not metric measurements). What survives
unconditionally is the categorical bridge design (below); the
quantitative pLGH-convergence reading is the repair target.

\medskip

\noindent\emph{Honest scope.} The bridge is K$^{+}$-weak-form only.
The strong-form bridge extending $\Wfun$ to the Paper~46~\cite{paper46}
Appendix~B enlarged substrate with chirality-flipping generators (Q1') is
deferred as a multi-month NCG-mathematics follow-on. The general
$\mathrm{G2}$ metric-level closure on non-wedge carriers remains the
multi-month frontier; we close only the wedge-specific case via the
bridge image on the synthetic side. The construction preserves
$\Wfun$'s functorial (not isometric) character throughout:\ the Krein
metametric $\BigDeth$ and the synthetic time separation $\ell^{L}$
measure different physical quantities (thermal vs.\ geometric time) of
the same wedge KMS state.

\medskip

\noindent\emph{Place in the math.OA standalone series.} Tenth in the
GeoVac math.OA arc (siblings:\ Papers
29~\cite{paper29}, 32~\cite{paper32}, 38~\cite{paper38},
39~\cite{paper39}, 40~\cite{paper40_unified}, 42~\cite{paper42},
43~\cite{paper43}, 44~\cite{paper44}, 45~\cite{paper45},
46~\cite{paper46}, 47~\cite{paper47}). The result composes
Paper~45/46 verbatim on the Krein-algebraic side with the
Mondino--S\"amann synthetic-Lorentzian Gromov--Hausdorff framework on
the synthetic side, mediated by the Connes--Rovelli thermal-time
identification.
\end{abstract}

\subsection*{Status and scope (June 2026)}
Three issues are established upstream of this paper: (i) the
$\Kplus$-restricted compressed-Dirac Lipschitz seminorm underlying
$\Lprop$ vanishes identically on the truncated operator system ---
Krein-self-adjointness of the anti-Hermitian spatial part
$i\,\DGV \otimes I$ forces $\{\JL, \DGV \otimes I\} = 0$, so the
$\Kplus$ compression annihilates the spatial Dirac, and the
momentum-diagonal temporal multipliers commute with $\partial_t$
(verified bit-exact at $(\nmax, \Nt) \in \{(2,3), (3,5)\}$); (ii)
the ``Latr\'emoli\`ere Thm~5.5 / Def.~3.4 / Def.~3.5'' references
cited for the tunneling-pair bound do not exist in the cited source
--- the device actually used is van~Suijlekom's state-space
Gromov--Hausdorff framework; (iii) the constant $2/(\nmax + 1)$ is
the maximum of the Plancherel mass distribution, not a cb-norm of
any map used in the proofs (the smoothing map is unital completely
positive, with cb-norm $1$).

Consequently the bridge functor's domain category $\KreinMM$, as
instantiated on the Paper~45 substrate, is \emph{degenerate as a
metric object}: the Lipschitz seminorms feeding the Krein metametric
have kernels far exceeding the scalars, so the instantiated objects
are not quantum (proper) metric spaces in the cited frameworks.
T3 (Theorem~\ref{thm:bit_exact_panel}) transports formula values,
not metric data; T6 (Theorem~\ref{thm:g2_wedge_closure}, the G2
metric-level closure for the GeoVac wedge) rests on the same
degenerate quantity and is \emph{descoped}. What survives is the
categorical design --- the Wick-rotation functor $\Wfun$ and the
Connes--Rovelli thermal-time reading resolving the F2
forward-vs-reverse triangle mismatch --- as a proposal conditional
on repair of the underlying substrate. The Lorentzian quantum-metric
convergence question is an open named research target.

\tableofcontents

% =====================================================================
\section{Introduction}\label{sec:intro}
% =====================================================================

\subsection{Motivation:\ the F2 forward-vs-reverse triangle mismatch}
\label{sec:motivation}

Two recent advances in the convergence theory of generalized Lorentzian
geometric structures, both published in December 2025, provide
complementary but \emph{a priori} incompatible frameworks:

\begin{itemize}
\item \textbf{Latr\'emoli\`ere's pinned proper quantum metric space
framework}~\cite{latremoliere2025_hypertopology} (arXiv:2512.03573,
December~3, 2025) extends the Latr\'emoli\`ere
propinquity~\cite{latremoliere2018, latremoliere_metric_st_2017} from
quantum \emph{compact} metric spaces to a quantum Gromov--Hausdorff
\emph{hypertopology} on the class of pointed proper quantum metric
spaces. The resulting metametric $\BigDeth$ satisfies a 4-point
\emph{relaxed forward} triangle inequality with codomain $[0, \infty]$.
The framework is operator-algebraic:\ objects are C*-algebras with
Lipschitz seminorms, topographies (Abelian C*-subalgebras), and pin
states; morphisms are topographic quantum M-isometries.

\item \textbf{Mondino--S\"amann's covered Lorentzian pre-length space
framework}~\cite{mondino_samann2025_pointed} (arXiv:2504.10380~v4,
December~9, 2025) extends the synthetic Lorentzian geometry of
Kunzinger--S\"amann pre-length spaces~\cite{kunzinger_samann2018} to a
pointed Lorentzian Gromov--Hausdorff (pLGH) convergence theory on
covered LPLS. The synthetic time separation $\ell$ satisfies the
3-point \emph{reverse} triangle inequality (causal-curve
super-additivity of proper time) with codomain
$\{-\infty\}\cup[0,\infty]$. The framework is synthetic-geometric:\
objects are sets with timelike-causal pre-orders, cover slabs, and
basepoint events.
\end{itemize}

The two frameworks share many structural elements:\ pinning, exhaustive
covers, slab bounds, ε-nets, pre-compactness, Gromov--Hausdorff
convergence. They are aimed at different aspects of the same physics:\
generalized-geometric approximation of Lorentzian objects. However,
the \emph{algebraic shapes} of their central inequalities are
incompatible:

\begin{center}
\begin{tabular}{lcc}
\toprule
& Latr\'emoli\`ere 2512.03573 & Mondino--S\"amann 2504.10380 \\
\midrule
codomain         & $[0, \infty]$                 & $\{-\infty\}\cup [0, \infty]$ \\
triangle direction & forward                     & reverse \\
triangle structure & 4-point relaxed              & 3-point \\
sign convention    & state-distance ($\ge 0$)     & proper time ($\ge 0$ when causal) \\
\bottomrule
\end{tabular}
\end{center}

We call this the \textbf{F2 mismatch}. Naive attempts at a bridge fail:
negating the time variable (R1) puts the codomain in the wrong
half-line and the 4-point structure does not collapse to 3-point under
sign reversal. The substantive question is whether a bridge can be
\emph{categorical}:\ a functor sending Krein-algebraic objects (where
relevant) to synthetic-Lorentzian objects (where relevant), with the
F2 mismatch interpreted as a structural signature of two distinct
projection mechanisms applied to the same physical object, rather than
as a defect to be reconciled.

\subsection{Main result:\ the Wick-rotation functor and the
Krein--Mondino--S\"amann Bridge}
\label{sec:main_result}

\paragraph*{What this paper contributes and what it inherits.}
From Mondino--S\"amann~\cite{mondino_samann2025_pointed} we inherit
the covered Lorentzian pre-length space framework, the synthetic time
separation $\ell$, the 3-point reverse triangle inequality, the
pLGH-convergence axioms (Def.~3.6), and the synthetic pre-compactness
theorem (Thm.~6.2). From
Latr\'emoli\`ere~\cite{latremoliere2025_hypertopology} we inherit the
pinned proper quantum metric space framework, the 4-point relaxed
forward metametric $\BigDeth$, and the hypertopology of admissible
truncated covers. From Connes--Rovelli~\cite{connes_rovelli1994} we
inherit the thermal-time hypothesis identifying modular flow with a
physical time parameter on a KMS state. \emph{What this paper adds}
is (i) the categorical bridge $\Wfun: \KreinMM \to \LorPLG$ between
those constructions; (ii) the F2 resolution via the
thermal-vs-geometric-time distinction (the forward/reverse mismatch
is a structural signature of two projection mechanisms applied to one
KMS state, not a defect of either framework); (iii) the Bridge
Theorem (B1)--(B4) at theorem-grade rigor on the K$^{+}$-weak-form
substrate; and (iv) the seven newly accessible theorems on the
GeoVac wedge image (T1--T6a), headlined by T6's wedge-specific
metric-level \emph{route toward} the published-open G2 problem of
Paper~45~\cite{paper45} \S~1.4 (descoped --- conditional on the descoped
Krein seminorm; the synthetic-side transport is a proposed route, not a closure). The strong-form bridge on the
Paper~46~\cite{paper46} Appendix~B enlarged substrate (Q1') is
explicitly deferred (Section~\ref{sec:open_q1prime}); we do
\emph{not} claim that or the general non-wedge G2 metric-level closure.

\medskip

We construct (Definition~\ref{def:wick_functor},
Theorem~\ref{thm:bridge}) the Wick-rotation functor
\[
\Wfun: \KreinMM \longrightarrow \LorPLG,
\]
contravariant by Gelfand duality, sending a Krein-pointed proper QMS
$(\Acal^{K}, L^{K}, \Mcal^{L}, \BWvac)$ to the covered LPLS
\[
\Wfun(\Acal^{K}, L^{K}, \Mcal^{L}, \BWvac)
\;=\; (\hat{\Mcal}^{L}, \;\ell^{L}, \;\hat{\omega}_{W}^{L}, \;\hat{\Ucal}^{L}),
\]
where $\hat{\Mcal}^{L}$ is the Gelfand spectrum of the M-diagonal
Abelian topography, $\ell^{L} = \kappag\cdot\tmod^{\BWvac}$ is
$\kappag$ times the modular-flow time on the wedge boost orbits
($-\infty$ off-orbit), $\hat{\omega}_{W}^{L}$ is the
Bisognano--Wichmann vacuum as character of the topography, and
$\hat{\Ucal}^{L}$ is the truncation cover.

The Krein--Mondino--S\"amann Bridge Theorem (Theorem~\ref{thm:bridge})
states that the \emph{structural} bridge properties (B1, B3) hold at
theorem-grade rigor, while the \emph{metric-level convergence} leg (B4)
is \textbf{descoped} --- it rests on the degenerate Krein metric of
Paper~45 (the Status note), so it is a conditional transport pending
substrate repair, not an established metric convergence:
\begin{itemize}
\item[(B1)] \emph{Structural correspondence}\ — $15$ rows of structural
alignment (set, topology, time separation, basepoint, cover, scales,
slabs, ε-nets, convergence, pre-compactness, etc.).
\item[(B2)] \emph{Reverse triangle inequality} on the bridge image\ —
on-orbit via additivity of Tomita--Takesaki modular flow at integer
$\Kalpha$ spectrum~\cite{paper42}; off-orbit \emph{structurally
vacuous} at the K$^{+}$-weak-form bridge because the M-diagonal
topography forces orbit-pair contradiction (Decomposition O Case (iii)
is empty).
\item[(B3)] \emph{Pre-compactness inheritance}\ — Mondino--S\"amann
Thm~6.2 applies on the bridge image via the Paper~44 propagation-number
$=2$ cardinality bound.
\item[(B4)] \emph{Convergence transport} (\textbf{descoped})\ — Krein-side
convergence under the Latr\'emoli\`ere hypertopology \emph{would} induce
Mondino--S\"amann pLGH convergence, with Plancherel-nested Berezin maps
matching the MS Def~3.6 axioms one-to-one; but the premise $\BigDeth\to0$
is built from the degenerate Krein metric of Paper~45, so this transport
is conditional pending substrate repair, not an established metric
convergence.
\end{itemize}

The F2 mismatch is resolved via the Connes--Rovelli thermal-time
hypothesis~\cite{connes_rovelli1994}:\ the Krein metametric
$\BigDeth$ measures \emph{thermal time} of the wedge KMS state, the
synthetic time separation $\ell^{L}$ measures \emph{geometric proper
time}, and the Wick rotation between them is the analytic continuation
mapping additive thermal-time structure to super-additive
geometric-time structure. Concretely, on the
Bisognano--Wichmann wedge,
$t_{\mathrm{geometric}} = \kappag \cdot t_{\mathrm{thermal}}$
by the four-witness Wick-rotation theorem of Paper~42~\cite{paper42}.

\subsection{GeoVac wedge application:\ seven newly accessible theorems}
\label{sec:wedge_headline}

The bridge is illustrated on the canonical Paper~43~\cite{paper43}
hemispheric wedge $W_{L} = P_{W}^{\spat} \otimes P_{t \ge 0}$ on the
compact-temporal Lorentzian Krein space $\Krein_{\nmax, \Nt, T}$, with
Bisognano--Wichmann vacuum $\BWvac = e^{-\Kalpha}/Z$ and modular
Hamiltonian $\Kalpha = J_{\mathrm{polar}}^{W}$ of integer spectrum
$\mathrm{two\_m}_{j}$. All five Krein PPQMS axioms verify by direct
construction. Seven newly accessible theorems follow
(Section~\ref{sec:wedge_application}):

\begin{itemize}
\item \textbf{T1} (wedge as MS-covered LPLS) — every finite cutoff
produces a well-defined Mondino--S\"amann covered Lorentzian pre-length
space with uniform timelike diameter $2\pi$ across all cover levels;
\item \textbf{T2} (synthetic compactness of wedge family) — pLGH
pre-compactness for sequences of GeoVac wedge cutoffs;
\item \textbf{T3} (synthetic bit-exact panel) — the synthetic pLGH-rate
panel at $(\nmax,\Nt)\in\{(2,3),(3,5),(4,7)\}$ inherits bit-identically
the Paper~45 Krein-side $\Lprop$ values $\{2.0746, 1.6101, 1.3223\}$
(conditional on the bridge;\ these are the descoped Paper~45
\emph{rate-formula} evaluations, not metric-convergence measurements ---
the Lorentzian quantum-metric convergence is descoped);
\item \textbf{T4} (synthetic four-witness readings) — the
Paper~42~\cite{paper42} four-witness Wick-rotation theorem's six
physical witnesses reduce to two synthetic-side equivalence classes
distinguished by the rapidity-rate normalization $\kappag$;
\item \textbf{T5} (synthetic Pythagorean orthogonality with M1
signature) — the operator-algebraic Pythagorean orthogonality
$\langle \Hloc, \Dwedge \rangle_{HS} = 0$ of Paper~43~\cite{paper43}
\S~10.2 transports to a synthetic-side foliation orthogonality, with
the closed-form $1/\pi^{2}$ residual prefactor identified as the
master Mellin engine M1 Hopf-base-measure signature
(Paper~32~\cite{paper32} \S~VIII);
\item \textbf{T6} (proposed \textbf{G2 metric-equivalent route for the
GeoVac wedge}, conditional/descoped) — the deepest design content:\ the
de-compactification limit $T \to \infty$ \emph{would} close at the
pLGH-equivalent (metric-level-equivalent) level on the synthetic side
for the GeoVac wedge, \emph{conditional on the bridge} --- which inherits
the descoped Paper~45 degeneracy. It is therefore a \emph{proposed}
workaround for the published-open non-compact
propinquity problem (Paper~45~\cite{paper45} \S~1.4 $\mathrm{G2}$) at
the GeoVac-wedge-specific level, not a closure;
\item \textbf{T6a} (synthetic three-carrier identification) —
Paper~47~\cite{paper47} \S~6 three-carrier identification (periodic /
Dirichlet / non-compact) extends from operator-algebraic to
synthetic-Lorentzian level.
\end{itemize}

\subsection{Honest scope}
\label{sec:honest_scope}

The bridge $\Wfun$ is constructed at the \textbf{K$^{+}$-weak-form}
level only (the substrate of Paper~45~\cite{paper45}). The strong-form
extension of $\Wfun$ to the Paper~46~\cite{paper46} Appendix~B
enlarged substrate (chirality-flipping generators $M^{\mathrm{flip}}$
with $\{J, M^{\mathrm{flip}}\} = 0$) is \textbf{not} constructed here;
it is the named open question \textbf{Q1'} (Section~\ref{sec:open_q1prime}).
On the enlarged substrate the off-orbit super-additivity ceases to be
vacuous and becomes a load-bearing operator-algebraic statement,
turning Q1' into a substantive multi-month NCG-mathematics target.

The functor $\Wfun$ is \textbf{functorial, not isometric}:\
$\BigDeth$ on the Krein-algebraic side and $\ell^{L}$ on the synthetic
side measure different physical quantities of the same wedge KMS
state. The forward-vs-reverse triangle direction is not a "mismatch
to be reconciled" but a structural signature of two distinct
projection mechanisms applied to the same physical object. This is
the F2 resolution.

For the GeoVac wedge, the substantive \textbf{T6} \emph{route toward}
$\mathrm{G2}$ is wedge-specific:\ it proposes (conditionally; descoped, as
it inherits the degenerate Krein seminorm) a de-compactification
limit at the propinquity-equivalent (pLGH) level on the synthetic side
for the wedge specifically. The general non-compact Latr\'emoli\`ere
propinquity problem (Paper~45~\cite{paper45} \S~1.4 $\mathrm{G2}$ on
arbitrary non-wedge carriers) remains the multi-month NCG-mathematics
frontier and is \textbf{not} closed by this paper. The Paper~47
\cite{paper47} norm-resolvent closure at the spectral level remains
the operator-algebraic gold standard for general carriers; T6
specializes the metric-level question to the wedge via the bridge.

\subsection{Roadmap of the paper}
\label{sec:roadmap}

Section~\ref{sec:setup} recalls the necessary background:\
Latr\'emoli\`ere's pinned proper QMS framework
(\S~\ref{sec:latremoliere_recap}), the Krein operator-system substrate
of Papers~44/45/46 (\S~\ref{sec:krein_substrate_recap}), the
Mondino--S\"amann pre-length space framework
(\S~\ref{sec:mondino_samann_recap}), the Bisognano--Wichmann vacuum
and the four-witness Wick-rotation theorem of Paper~42
(\S~\ref{sec:bw_four_witness_recap}), and the Connes--Rovelli
thermal-time hypothesis (\S~\ref{sec:connes_rovelli_recap}).

Section~\ref{sec:krein_ppqms} establishes the Krein-pointed proper
quantum metric space:\ a 4-tuple
$\mathbb{X}^{K} = (\Acal^{K}, L^{K}, \Mcal^{L}, \BWvac)$
in which the nine Latr\'emoli\`ere axioms transport at the
\emph{categorical/structural} level. \textbf{(Descoped.)} On the
Paper~45 substrate the Krein Lipschitz seminorm $L^{K}$ is degenerate
(its kernel exceeds the scalars; Status note (i)), so $\mathbb{X}^{K}$
is a categorical-design construction, \emph{not} an established proper
quantum metric space in the cited sense; the metric-level content of
\S~\ref{sec:krein_ppqms} is descoped accordingly.

Section~\ref{sec:krein_hypertopology} constructs the Krein
hypertopology on the class of Krein-pointed proper QMS, with explicit
M-tunnel extent, M-local metametric, 4-point relaxed triangle
inequality, coincidence theorem, and inframetric structure. The
compact-case agreement at $\Nt = 1$ reduces bit-exactly to the
Paper~38~\cite{paper38} SU(2) state-space Gromov--Hausdorff hypertopology.

Section~\ref{sec:bridge} constructs the Wick-rotation functor and
states the Bridge Theorem. The F2 mismatch is resolved via R3
(Connes--Rovelli thermal time) as the unifying framework. The
candidate resolutions R1 (negate time) and R2 (sub-functors) are
tested and found insufficient.

Section~\ref{sec:wedge_application} applies the bridge to the GeoVac
hemispheric wedge, with the seven newly accessible theorems.

Section~\ref{sec:compactness} states the synthetic-Lorentzian
compactness theorems, with T6 (G2 metric-level closure for the GeoVac
wedge) as the headline.

Section~\ref{sec:open_questions} catalogues the open questions:\ Q1'
strong-form bridge with enlarged substrate, G2 on the general
non-wedge carrier, G3 cross-manifold, G4 inner-factor calibration.

\subsection{Related work}
\label{sec:related_work}

The Latr\'emoli\`ere lineage~\cite{latremoliere2018,
latremoliere_metric_st_2017, latremoliere_metric_propinquity_2015,
latremoliere2025_hypertopology, latremoliere_pp_2018af,
farsi_latremoliere2024, farsi_latremoliere2025} provides the
operator-algebraic side of the bridge. The pre-2025 papers handle the
compact case; the 2512.03573 hypertopology paper~\cite{latremoliere2025_hypertopology}
is the first to extend to proper (locally compact, non-unital) quantum
metric spaces. Our Krein lift inherits and extends this framework.

The Mondino--S\"amann lineage~\cite{mondino_samann2025_pointed,
kunzinger_samann2018, mondino_ryborz_samann2025} provides the
synthetic-Lorentzian side. The 2504.10380 v4 paper~\cite{mondino_samann2025_pointed}
introduces the pointed extension with covered LPLS and pLGH
convergence; earlier papers by Kunzinger--S\"amann established the
underlying Lorentzian pre-length-space framework. The
Mondino--Ryborz--S\"amann stability paper~\cite{mondino_ryborz_samann2025}
(May 2026) is the most recent flagship in this lineage.
A concurrent extension by Braun--S\"amann~\cite{braun_samann2025}
develops a Gromov reconstruction theorem and measured Gromov--Hausdorff
convergence in Lorentzian geometry; the bridge constructed here targets
the unmeasured pointed framework of~\cite{mondino_samann2025_pointed},
but the measured extension is a natural future bridge target.

Adjacent synthetic-Lorentzian convergence frameworks include:\
Sormani--Vega null distance~\cite{sormani_vega2016},
Sakovich--Sormani timed Gromov--Hausdorff~\cite{sakovich_sormani2024},
Minguzzi--Suhr bounded Lorentzian metric
spaces~\cite{minguzzi_suhr2024}, M\"uller's Cauchy
slabs~\cite{muller2022}, Allen--Burtscher metric
GH~\cite{allen_burtscher2022}, Che--Perales--Sormani timed
Hausdorff~\cite{che_perales_sormani2025}. These frameworks operate in
distinct categories from the present bridge (Mondino--S\"amann's
ε-net-of-causal-diamonds approach versus null-distance metric-GH or
timed Hausdorff variants); they are cited for completeness but are not
in direct competition with the bridge construction.

The Connes--Rovelli thermal-time hypothesis~\cite{connes_rovelli1994}
supplies the unifying interpretation. Active research on
Connes-spectral-triples plus Tomita--Takesaki modular structure as a
NCG framework for general relativity (e.g., Bertozzini--Conti--Lewkeeratiyutkul
\cite{bertozzini_conti_lewkeeratiyutkul2009}) is adjacent to R3 but
does not target the specific Krein-pointed PPQMS $\to$
Mondino--S\"amann pLGH bridge constructed here.

Recent NCG-side work on twisted spectral triples and pseudo-Riemannian
spectral geometry includes Bizi--Brouder--Besnard~\cite{bizi_brouder_besnard2018},
van den Dungen~\cite{vandungen2016}, Franco--Eckstein~\cite{franco_eckstein2014},
Strohmaier~\cite{strohmaier2006}, and the recent Nieuviarts
twisted-spectral-triple paper~\cite{nieuviarts2025_v2} (December~2025
v2). The Nieuviarts framework is restricted to even KO-dimension and
does not apply directly to our KO-dim~3 substrate; we cite it as an
alternative-to-Wick-rotation route in the Lorentzian-extension
program. Ketterer~\cite{ketterer2026} (May 2026) and
Kubota~\cite{kubota2026} (May 2026) are recent synthetic-side
internal refinements; the former treats convergence of Lorentzian
spaces and curvature bounds for generalized cones, the latter
establishes a Lorentzian coarea inequality. We cite both as adjacent works in the synthetic-Lorentzian
convergence literature.

The Hekkelman--McDonald papers~\cite{hekkelman_mcdonald2024} on Connes spectral triples and
noncommutative integration provide non-compact propinquity machinery
that does not extend to a metric-level closure on the
$\sthree \times \R_{t}$ carrier; our T6 closure for the wedge sits on
the synthetic side via the bridge rather than via direct non-compact
propinquity.

To the author's knowledge, no published paper constructs a categorical
bridge between operator-algebraic Lorentzian propinquity and synthetic
Lorentzian Gromov--Hausdorff frameworks. The bridge constructed here
is novel.

% =====================================================================
\section{Setup}\label{sec:setup}
% =====================================================================

\subsection{Latr\'emoli\`ere's pinned proper QMS
framework}\label{sec:latremoliere_recap}

We recall the central definitions from
Latr\'emoli\`ere~\cite{latremoliere2025_hypertopology}.

\begin{definition}[Latr\'emoli\`ere 2512.03573 Def~1.18:\
Leibniz seminorm]\label{def:leibniz}
Let $\Acal$ be a C*-algebra. A \emph{Leibniz Lipschitz seminorm}
$L$ on $\Acal$ is a (typically unbounded) seminorm on a dense
*-subalgebra $\mathrm{dom}(L) \subseteq \Acal$ satisfying $L(a^{*}) = L(a)$
and the Leibniz inequality
$L(ab) \le \norm{a} L(b) + L(a) \norm{b}$.
\end{definition}

\begin{definition}[Latr\'emoli\`ere 2512.03573 Def~1.22:\ separable
QLCMS]\label{def:sep_qlcms}
A \emph{separable quantum locally compact metric space} (QLCMS) is a
pair $(\Acal, L)$ of a separable C*-algebra $\Acal$ and a Leibniz
seminorm $L$ satisfying:
\begin{itemize}
\item[(FM)] Fortet--Mourier distance metrizes the weak-$*$
topology on $\Scal(\Acal)$;
\item[(CB)] the unit ball $\{a \in \mathrm{dom}(L) : L(a) \le 1\}$ is
closed;
\item[(B)] there exists a strictly positive $h \in \Acal_{\mathrm{sa}}$
and a state $\mu \in \Scal(\Acal)$ with $\mu(h) = 1$ such that
\[
\{h a h : a = b + t\mathbf{1},\; L(b) \le 1,\; \mu(b) = -t\}
\]
is bounded.
\end{itemize}
\end{definition}

\begin{definition}[Latr\'emoli\`ere 2512.03573 Def~1.26:\ pinned
QLCMS]\label{def:pinned_qlcms}
A \emph{pinned QLCMS} is a triple $(\Acal, L, \mu)$ extending
Definition~\ref{def:sep_qlcms} with $\mu \in \Scal(\Acal)$ a pin
state, such that $\{\phi \in \Scal(\Acal) : \mathrm{mk}_{L}(\phi, \mu)
< \infty\}$ is weak-$*$ dense in $\Scal(\Acal)$, where $\mathrm{mk}_{L}$
is the Monge--Kantorovich distance.
\end{definition}

\begin{lemma}[Latr\'emoli\`ere 2512.03573 Lemma~1.25]\label{lem:tightness_corollary}
If $(\Acal, L)$ is a separable QLCMS with the (B)-axiom witnessed by
$(h, \mu)$, then $\{\phi : \mathrm{mk}_{L}(\phi, \mu) < \infty\}$ is
weak-$*$ dense in $\Scal(\Acal)$. In particular, the (B)-axiom of
Definition~\ref{def:sep_qlcms} implies the density condition of
Definition~\ref{def:pinned_qlcms}.
\end{lemma}

\begin{definition}[Latr\'emoli\`ere 2512.03573 Def~1.37 + 1.42:\
pointed proper QMS]\label{def:pointed_proper_qms}
A \emph{pinned proper QMS} is a pinned QLCMS
$(\Acal, L, \mu)$ containing an $L$-Lipschitz $\mu$-pinned
exhaustive sequence $(h_{n}) \subseteq \mathrm{dom}(L)$ with
$L(h_{n}) \to 0$, $\mu(h_{n}) \to 1$, $\norm{h_{n}} \to 1$. A
\emph{topography} $\mathfrak{M} \subseteq \Acal$ is an Abelian
C*-subalgebra containing such an exhaustive sequence. A
\emph{pointed proper QMS} is a 4-tuple $(\Acal, L, \mathfrak{M}, \mu)$
with $\mathfrak{M}$ a topography and $\mu$ restricting to a character
of $\mathfrak{M}$.
\end{definition}

The associated structures are M-isometries, M-tunnels, M-local
metametric $\eth_{M}$, parameter-family metametric $\BigDeth$, and a
4-point relaxed forward triangle:
\[
\BigDeth(\mathbb{A}, \mathbb{B})
\le (1 + \BigDeth(\mathbb{A}, \mathbb{D}))^{2}\BigDeth(\mathbb{D}, \mathbb{B})
+ (1 + \BigDeth(\mathbb{D}, \mathbb{B}))^{2}\BigDeth(\mathbb{A}, \mathbb{D}).
\]
The Latr\'emoli\`ere coincidence theorem (Thm~3.6) characterizes
$\BigDeth = 0$ via the existence of a full topographic M-isometry.

\subsection{Krein operator-system substrate from
Papers~44/45/46}\label{sec:krein_substrate_recap}

The Krein-side substrate is the compact-temporal Lorentzian Krein
space
\[
\Krein_{\nmax, \Nt, T}
= \HGV^{\nmax} \otimes \C^{\Nt}
\]
of Paper~44~\cite{paper44}, with fundamental symmetry $\JL = J_{\spat}
\otimes I_{\Nt}$ where $J_{\spat}^{2} = +I$ is the spatial chirality
swap; the K$^{+}$-restricted Hilbert space is
$\Kplus = \{|\psi\rangle : \JL |\psi\rangle = +|\psi\rangle\}$.
The Lorentzian Dirac is (van den Dungen 2016 Prop.~4.1
\cite{vandungen2016})
\[
\DL = i\bigl(\gamma^{0}\otimes \partial_{t} + \DGV \otimes I_{\Nt}\bigr),
\]
Krein-self-adjoint:\ $\DL^{\times} = \DL$ where $\dagger$ is the
Krein adjoint $\gamma^{0}\cdot *\cdot\gamma^{0}$. The natural Lorentzian
operator system is
\[
\Op^{L} = \mathrm{span}_{\C}\{M^{\spat}_{N, L, M} \otimes M^{\temp}_{p}
:\ N \le \nmax,\ L < N,\ |M| \le L,\ 0 \le p \le \Nt - 1\},
\]
of finite dimension at finite cutoff, with propagation number
$\mathrm{prop}(\Op^{L}) = 2$ under the Weyl-doubled achievable envelope
(Paper~44~\cite{paper44} Prop.~prop=2).

Paper~45~\cite{paper45} is the K$^{+}$-restricted weak-form
\emph{degeneracy} theorem:\ the natural Lorentzian-propinquity seminorm
$\Lprop$ vanishes identically, so the would-be convergence
\[
\Lprop(\Tcal^{L}_{\nmax, \Nt, T}, \Tcal^{L}_{\Manifold})
\;\le\; \Cthreejoint \cdot \gammajoint_{\nmax, \Nt, T} \;\to\; 0
\]
is a closed-form \emph{rate-formula} evaluation (panel values
$\Lprop(2,3) = 2.0746$, $\Lprop(3,5) = 1.6101$, $\Lprop(4,7) = 1.3223$),
\emph{not} a quantum-metric distance --- the Lorentzian quantum-metric
convergence is descoped. Paper~46~\cite{paper46}
carries the same rate formula to a strong-form (no K$^{+}$ restriction)
on the natural chirality-doubled scalar-multiplier substrate
($\Lambda^{\mathrm{strong}} = \Lprop$, also degenerate), and proposes a
route on a further enlarged substrate (Appendix~B) with
chirality-flipping generators $M^{\mathrm{flip}}$ satisfying
$\{\JL, M^{\mathrm{flip}}\} = 0$ --- also descoped.

\subsection{Mondino--S\"amann pre-length space
framework}\label{sec:mondino_samann_recap}

We recall the central definitions from
Mondino--S\"amann~\cite{mondino_samann2025_pointed}.

\begin{definition}[Mondino--S\"amann Def~2.3:\ Lorentzian pre-length
space]\label{def:lpls}
A \emph{Lorentzian pre-length space} (LPLS) is a pair $(X, \ell)$
with $X$ a set carrying a topology finer than the chronological
topology, and $\ell: X \times X \to \{-\infty\} \cup [0, \infty]$
satisfying $\ell(x, x) \ge 0$ and the \emph{reverse triangle
inequality}
\begin{equation}\label{eq:reverse_triangle}
\ell(x, y) + \ell(y, z) \le \ell(x, z) \quad \forall x, y, z \in X
\end{equation}
(with convention $\pm\infty + \mp\infty = 0$ on the LHS yielding
$-\infty$ when undefined). Timelike, causal, and causal-diamond
relations are derived in the standard way:\ $x \ll y \iff \ell(x, y) > 0$,
$x \le y \iff \ell(x, y) \ge 0$, $J(x, y) = \{z : x \le z \le y\}$,
and the time separation $\tau := \max(0, \ell)$.
\end{definition}

\begin{definition}[Mondino--S\"amann Def~3.2:\ ε-net]\label{def:eps_net}
An \emph{ε-net} for $A \subseteq X$ is a collection
$\mathcal{S} = (J_{i})_{i \in \Omega}$ of causal diamonds with
$\tau(J_{i}) \le \varepsilon$ and $A \subseteq \bigcup_{i} J_{i}$.
\end{definition}

\begin{definition}[Mondino--S\"amann Def~3.6:\ LGH convergence of
subsets]\label{def:lgh_convergence}
A sequence $A_{n} \to A$ converges in the Lorentzian
Gromov--Hausdorff (LGH) sense if there exist ε-nets $S_{n}$ for
$A_{n}$ and $S$ for $A$ with:\
(i) matching cardinality;
(ii) correspondences $R_{n}$ between vertices with $\mathrm{dis}(R_{n})
\to 0$;
(iii) extension property of correspondences;
(iv) forward density.
\end{definition}

\begin{definition}[Mondino--S\"amann Def~3.8:\ covered LPLS]\label{def:covered_lpls}
A \emph{covered LPLS} is a 4-tuple $(X, \ell, o, \Ucal)$ where
$(X, \ell)$ is an LPLS, $o \in X$ is a basepoint event, and
$\Ucal = (U_{k})_{k \in \N}$ is a countable cover satisfying:
(i) $\bigcup_{k} U_{k} = X$;
(ii) $U_{k} \subseteq U_{k+1}$ (nested);
(iii) $o \in U_{k}$ for all $k$;
(iv) $\sup_{x, y \in U_{k}} \tau(x, y) < \infty$ (slab bounds).
\end{definition}

\begin{definition}[Mondino--S\"amann Def~3.12:\ pLGH
convergence]\label{def:plgh}
$(X_{n}, \ell_{n}, o_{n}, \Ucal_{n}) \to (X, \ell, o, \Ucal)$ in the
\emph{pointed LGH} sense iff for each $k$, $U_{k, n} \to U_{k, \infty}$
in LGH sense.
\end{definition}

The Mondino--S\"amann pre-compactness theorem (Thm~6.2) gives that
any sequence of covered LPLS with (i) uniform slab bounds, (ii) uniform
ε-net cardinality bounds, (iii) cumulative ε-net nesting has a
pLGH-convergent subsequence.

\subsection{Bisognano--Wichmann vacuum and the four-witness
Wick-rotation theorem}\label{sec:bw_four_witness_recap}

The Paper~43~\cite{paper43} hemispheric wedge is
\[
W_{L} := P_{W}^{\spat} \otimes P_{t \ge 0},
\quad
P_{W}^{\spat} = \tfrac{1}{2}(I + R_{\mathrm{polar}}),
\]
where $R_{\mathrm{polar}}$ is the spatial $m_{j}$-reflection involution
(Paper~42~\cite{paper42} \S~4) and $P_{t \ge 0}$ is the diagonal
projection on $\C^{\Nt}$ to positive temporal indices. The
Bisognano--Wichmann vacuum state under the BW-α geometric
Hamiltonian (with $\beta = 2\pi/\kappag$) is
\[
\rho_{W}^{L} := e^{-\Kalpha}/Z,
\quad
\Kalpha = J_{\mathrm{polar}}^{W} \text{ on the wedge},
\]
with integer spectrum $\mathrm{two\_m}_{j} \in \Z_{>0}$. The
Tomita--Takesaki modular Hamiltonian $K_{\mathrm{TT}} = -\log \Delta$
is constructed from the wedge KMS state via polar decomposition
(Paper~42~\cite{paper42} \S~6).

The Paper~42~\cite{paper42} four-witness Wick-rotation theorem states
that six physical witnesses (Bisognano--Wichmann canonical with
$\kappag = 1$, Hartle--Hawking at $M = 1, 2$, Sewell at $M = 1$, Unruh
at $a = 1, 2$) all yield bit-identical modular operators $\Delta$
and Hamiltonians $K_{\mathrm{TT}}$ under the BW unit normalization
$\Hloc := \Kalpha / \beta$. The Riemannian-side bit-exact identity
$\sigma_{2\pi}^{\alpha}(O) = O$ holds at every finite cutoff
$\nmax \in \{1, 2, 3, 4, 5\}$, and the Tomita BW-γ flow $\sigma_{t}^{TT}$
satisfies the conjugacy $\sigma_{t}^{TT}(a) = \sigma_{-t}^{\alpha}(a)$
bit-exact. Paper~43~\cite{paper43} extends to Lorentzian signature via
BBB Krein construction~\cite{bizi_brouder_besnard2018} and verifies
the same six-witness collapse at every $(\nmax, \Nt)$ panel cell.

\subsection{Connes--Rovelli thermal-time
hypothesis}\label{sec:connes_rovelli_recap}

Connes and Rovelli~\cite{connes_rovelli1994} (1994) proposed:

\medskip

\noindent\textbf{Thermal-time hypothesis.} \emph{In quantum
statistical mechanics on a generally covariant system, the natural
physical time evolution at thermal equilibrium with state $\omega$ is
given by the modular automorphism group $\sigma_{t}^{\omega}$ of
$\omega$ (a state on the observable algebra $\Acal$).}

\medskip

Applied to a KMS state at inverse temperature $\beta$ on a Type~III
von Neumann algebra arising as the local algebra of a region of
spacetime (Bisognano--Wichmann setting), the thermal time of $\omega$
relates to the geometric proper time of the underlying spacetime by
\begin{equation}\label{eq:connes_rovelli_identification}
t_{\mathrm{thermal}} = \frac{\beta}{2\pi} \cdot t_{\mathrm{geometric}}
= \frac{1}{\kappag} \cdot t_{\mathrm{geometric}},
\end{equation}
where $\kappag$ is the surface gravity (equivalently the BW
Hawking-temperature parameter). This is the four-witness
Wick-rotation theorem of Paper~42~\cite{paper42} at the kinematic
level:\ modular flow IS Lorentz-boost orbit on the BW wedge, with
proportionality factor $1/\kappag = \beta/(2\pi)$.

\subsection{Notation and conventions}\label{sec:notation}

Throughout the paper:\
$\Krein = \Krein_{\nmax, \Nt, T} = \HGV^{\nmax} \otimes \C^{\Nt}$
denotes the compact-temporal Krein space;
$\JL = J_{\spat} \otimes I_{\Nt}$ the fundamental symmetry;
$\Kplus, \Kminus$ the K$^{\pm}$-eigenspaces of $\JL$;
$\Op^{L}$ the natural Lorentzian operator system;
$\DL$ the Krein-self-adjoint Lorentzian Dirac;
$\BWvac = e^{-\Kalpha}/Z$ the BW vacuum;
$\Kalpha = J_{\mathrm{polar}}^{W}$ the BW geometric Hamiltonian (with
integer spectrum $\mathrm{two\_m}_{j}$);
$\sigma_{t}^{\BWvac}$ the Tomita--Takesaki modular flow of $\BWvac$;
$\tmod^{\BWvac}(\hat{\omega}, \hat{\omega}') := \inf\{t \ge 0 :
\hat{\omega}' = \hat{\omega} \circ \sigma_{t}^{\BWvac}\}$ the
modular-flow time.

We write $\mathrm{Spec}(\mathfrak{B})$ for the Gelfand spectrum of an
Abelian C*-algebra $\mathfrak{B}$ and $\hat{\omega}$ for the character
induced by a state $\omega$ restricting to $\mathfrak{B}$.

% =====================================================================
\section{Krein-pointed proper quantum metric spaces}\label{sec:krein_ppqms}
% =====================================================================

We lift Latr\'emoli\`ere's pinned proper QMS framework to the
Krein-algebraic setting. Throughout this section, the substrate is
the K$^{+}$-restricted Hilbert space; the natural-substrate generators
all commute with $\JL$, so the Krein operator norm coincides with the
Hilbert operator norm on the K$^{+}$ restriction.
\emph{Descope note (read all of this section under it):} on the
Paper~45 substrate the Krein Lipschitz seminorm $L^{K}$ is degenerate
(its kernel exceeds the scalars; Status note (i)), so the axiom
verifications below establish a \emph{categorical/structural} lift, not
a metric (proper quantum metric space) in the cited sense; the
metric-level reading is descoped pending repair.

\subsection{The Krein-Leibniz Lipschitz seminorm}\label{sec:krein_leibniz}

\begin{definition}[Krein-Leibniz Lipschitz seminorm]\label{def:krein_leibniz}
On the natural Lorentzian operator system $\Op^{L}$ (or its continuum
limit on $\Tcal^{L}_{\Manifold}$), define
\[
L^{K}(a) := \opnorm{[\DL, a]}, \qquad a \in \mathrm{dom}(L^{K}) \subseteq \Op^{L},
\]
where $\opnorm{\cdot}$ is the operator norm on $\Bcal(\Krein)$
restricted to $\Kplus$. We call $L^{K}$ the \emph{Krein-Leibniz
Lipschitz seminorm}.
\end{definition}

\begin{lemma}[Krein-Leibniz inequality]\label{lem:krein_leibniz}
For every $a, b$ in the natural chirality-doubled scalar-multiplier
substrate,
\[
L^{K}(ab) \le \norm{a}_{\Op^{L}} L^{K}(b) + L^{K}(a) \norm{b}_{\Op^{L}}.
\]
\end{lemma}

\begin{proof}
On the natural substrate, all elements of $\Op^{L}$ commute with
$\JL$~\cite{paper44}, so both $a, b$ restrict cleanly to $\Kplus$.
On $\Kplus$, the standard Hilbert operator norm
$\opnorm{\cdot}|_{\Kplus}$ agrees with the algebra norm
$\norm{\cdot}_{\Op^{L}}$. The standard operator-algebra Leibniz
identity
$[\DL, ab] = [\DL, a] b + a [\DL, b]$ holds in $\Bcal(\Krein)$;
restricting to $\Kplus$ and applying submultiplicativity yields the
Krein-Leibniz inequality.
\end{proof}

\begin{remark}[Krein vs Hilbert operator norm on the enlarged
substrate]\label{rem:krein_op_norm}
The natural-substrate version of Lemma~\ref{lem:krein_leibniz} uses
the Hilbert operator norm on $\Kplus$. On the Paper~46~\cite{paper46}
Appendix~B \emph{enlarged} substrate with chirality-flipping
generators $M^{\mathrm{flip}}$ satisfying $\{\JL, M^{\mathrm{flip}}\} = 0$,
the K$^{+}$ restriction is no longer well-defined as a multiplier
(such generators map $\Kplus$ to $\Kminus$), and the two natural
operator norms (Krein operator norm with respect to the indefinite
Krein product vs.\ Hilbert operator norm on the full Krein space)
differ. The K$^{+}$-weak-form route of this paper makes the
distinction invisible; the strong-form route makes it visible at the
level of the gradient-norm absorption mechanism of
Paper~46~\cite{paper46} Appendix~B. For Q1' (Section~\ref{sec:open_q1prime})
this distinction becomes load-bearing.
\end{remark}

\subsection{The Krein-separable QLCMS and Krein-pinned QLCMS}\label{sec:krein_sep_pinned}

\begin{definition}[Krein-separable QLCMS]\label{def:krein_sep_qlcms}
A \emph{Krein-separable QLCMS} is a pair $(\Acal^{K}, L^{K})$ where
$\Acal^{K}$ is the natural Lorentzian operator system on $\Krein$
(at finite cutoff:\ the C*-algebra generated by $\Op^{L}$ in
$\Bcal(\Kplus)$, a finite-dimensional matrix algebra; in the
continuum:\ $\Acal^{K} = C(\Tcal^{L}_{\Manifold})$ acting on
$\Kplus_{\mathrm{continuum}}$), and $L^{K}$ is the Krein-Leibniz
Lipschitz seminorm on a dense *-subalgebra
$\mathrm{dom}(L^{K}) \subseteq \Acal^{K}$, satisfying:
\begin{itemize}
\item[(FM-K)] The Krein-positive Fortet--Mourier distance metrizes
the weak-$*$ topology on $\Scal(\Acal^{K})$;
\item[(CB-K)] The unit ball $\{a \in \mathrm{dom}(L^{K}) : L^{K}(a)
\le 1\}$ is closed in the operator-norm topology;
\item[(B-K)] There exists strictly positive $h^{L} \in \Acal^{K}_{\mathrm{sa}}$
and a state $\mu^{L} \in \Scal(\Acal^{K})$ with $\mu^{L}(h^{L}) = 1$
such that
\[
\{h^{L} a h^{L} : a = b + t\mathbf{1},\;
L^{K}(b) \le 1,\; \mu^{L}(b) = -t\}
\]
is bounded.
\end{itemize}
\end{definition}

\begin{lemma}\label{lem:krein_sep_qlcms_verification}
The pair $(\Acal^{K}, L^{K})$ with the natural Lorentzian operator
system substrate is a Krein-separable QLCMS, with witnessing
$(h^{L}, \mu^{L}) = (\Kalpha / Z, \BWvac)$.
\end{lemma}

\begin{proof}[Proof sketch]
The (FM-K), (CB-K), (B-K) axioms transport from Papers~42--45 substrate:
(FM-K) via Paper~44~\cite{paper44} \S~6 Krein-positive
Wasserstein--Kantorovich SDP framework, with weak-$*$ metrization
following Latr\'emoli\`ere~\cite{latremoliere2025_hypertopology}
Thm~1.9; (CB-K) via Banach--Alaoglu plus lower-semicontinuity of
$L^{K}$; (B-K) via Paper~42~\cite{paper42} \S~5 integer spectrum of
$\Kalpha$ ($h^{L}$ positive definite, $\norm{h^{L}} = 1$, $\BWvac(h^{L}) = 1$
by BW pinning) combined with Paper~44~\cite{paper44} propagation
number $= 2$ (bounded sets in finite-dimensional subspaces are bounded).
See A.2'~\S~2.3 of the synthesis memo for full detail.
\end{proof}

\begin{definition}[Krein-pinned QLCMS]\label{def:krein_pinned_qlcms}
A \emph{Krein-pinned QLCMS} is a triple $(\Acal^{K}, L^{K}, \BWvac)$
with $\BWvac \in \Scal(\Acal^{K})$ the BW vacuum, such that
$\{\phi \in \Scal(\Acal^{K}) : \mathrm{mk}_{L^{K}}(\phi, \BWvac) < \infty\}$
is weak-$*$ dense in $\Scal(\Acal^{K})$.
\end{definition}

\begin{lemma}[Continuum tightness via
Lemma~\ref{lem:tightness_corollary}]\label{lem:krein_continuum_tightness}
$(\Acal^{K}, L^{K}, \BWvac)$ is a Krein-pinned QLCMS. Equivalently:\
the density condition of Definition~\ref{def:krein_pinned_qlcms} holds.
\end{lemma}

\begin{proof}
By Lemma~\ref{lem:krein_sep_qlcms_verification}, the (B-K) axiom holds
with $(h^{L}, \mu^{L}) = (\Kalpha/Z, \BWvac)$. By
Lemma~\ref{lem:tightness_corollary} (Latr\'emoli\`ere
2512.03573 Lemma~1.25), the density of finite-Monge--Kantorovich
states follows mechanically.
\end{proof}

\begin{remark}\label{rem:tightness_dissolves}
Lemma~\ref{lem:krein_continuum_tightness} dissolves a caveat that
the diagnostic phase of the program had flagged as a 2--3 week
sub-sprint:\ the continuum tightness condition is a \emph{corollary}
of the (B-K) boundedness axiom via Latr\'emoli\`ere's own
Lemma~1.25. The Krein lift requires no separate verification beyond
the (B-K) ingredients.
\end{remark}

\subsection{The Krein L-Lipschitz $\BWvac$-pinned exhaustive sequence
and the approximate-unit theorem}\label{sec:krein_exhaustive_approx_unit}

\begin{definition}[Krein L-Lipschitz $\BWvac$-pinned exhaustive
sequence]\label{def:krein_exhaustive}
A \emph{Krein L-Lipschitz $\BWvac$-pinned exhaustive sequence}
is a sequence $(h_{n})_{n \in \N} \subseteq \mathrm{dom}_{\mathrm{sa}}(L^{K})$
with:
\begin{enumerate}
\item $\lim_{n\to\infty} L^{K}(h_{n}) = 0$,
\item $\lim_{n\to\infty} \BWvac(h_{n}) = 1$,
\item $\lim_{n\to\infty} \norm{h_{n}}_{\Acal^{K}} = 1$.
\end{enumerate}
\end{definition}

\begin{theorem}[Krein approximate-unit, transport of
Latr\'emoli\`ere 2512.03573 Thm~1.30]\label{thm:krein_approx_unit}
If $(\Acal^{K}, L^{K}, \BWvac)$ is a Krein-pinned QLCMS and
$(h_{n})$ is a Krein L-Lipschitz $\BWvac$-pinned exhaustive sequence,
then $(h_{n})$ is an approximate unit for $\Acal^{K}$.
\end{theorem}

\begin{proof}[Proof sketch]
Direct transport of Latr\'emoli\`ere 2512.03573 Thm~1.30 via Claims
1.31--1.34:\ weak-$*$ convergence to $1$ (using
Lemma~\ref{lem:krein_continuum_tightness}); $h_{n}^{2}$ pinning (via
Lemma~\ref{lem:krein_leibniz} Leibniz applied to $h_{n}\cdot h_{n}$);
$(1 - h_{n}^{2})$ vanishing on compact sets of states (via
Arzel\`a--Ascoli on the Fortet--Mourier metric ball); convergence to
$1$ in strict topology (via Kadison functional calculus on the
K$^{+}$-restricted commutative envelope of $h_{n}$).
\end{proof}

\textbf{The explicit Krein-side exhaustive sequence.} The natural
Krein-side exhaustive sequence is the truncated Krein triple
projector sequence:\
\[
h_{n} := P_{\nmax(n), \Nt(n), T(n)},
\]
the orthogonal projection onto the $\HGV^{\nmax(n)} \otimes \C^{\Nt(n)}$
truncation, with $(\nmax(n), \Nt(n), T(n)) \to (\infty, \infty, T_{\infty})$
along admissible scaling per Paper~47~\cite{paper47}. Verification:\
(i) $L^{K}(h_{n}) \le \gammajoint_{\nmax(n), \Nt(n), T(n)} \cdot O(1) \to 0$
by Paper~45~\cite{paper45} Sub-Sprint~C approximate-identity property;
(ii) $\BWvac(h_{n}) \to 1$ by support of $\BWvac$ on the wedge;
(iii) $\norm{h_{n}} = 1$ since $h_{n}$ is a projector.

\subsection{Krein topography and Krein-pointed proper
QMS}\label{sec:krein_topography}

\begin{definition}[Krein topography]\label{def:krein_topography}
A \emph{Krein topography} $\Mcal^{L}$ of a Krein-separable QLCMS
$(\Acal^{K}, L^{K})$ is an Abelian C*-subalgebra of $\Acal^{K}$
containing a strictly positive $h^{L} \in \mathrm{dom}(L^{K})$ and
admitting a character $\mu^{L}$ of $\Mcal^{L}$ witnessing the (B-K)
boundedness axiom restricted to $\Mcal^{L}$.
\end{definition}

\begin{lemma}[The natural Krein topography]\label{lem:natural_krein_topography}
The natural Krein topography on the GeoVac substrate is the
sub-operator-system
\[
\Mcal^{L} := \mathrm{span}_{\C}\bigl\{
M^{\spat}_{N, L, 0} \otimes I_{\Nt} :\ N \le \nmax,\ L < N
\bigr\}
\subseteq \Acal^{K},
\]
i.e.\ the chirality-doubled scalar multipliers with M-quantum-number
$M = 0$ tensored with the temporal identity.
\end{lemma}

\begin{proof}[Proof sketch]
\emph{Abelian:\ } Each $M^{\spat}_{N, L, 0}$ is a multiplication
operator by a spatial 3-Y function $Y^{(3)}_{N, L, 0}$, hence the
algebra is Abelian; tensoring with $I_{\Nt}$ preserves the structure.

\emph{Strictly positive $h^{L}$:\ } Take $h^{L} = \Kalpha/Z$; the BW
geometric Hamiltonian is M-diagonal and positive on the wedge by
Paper~42~\cite{paper42} \S~5.

\emph{(B-K) restricted to $\Mcal^{L}$:\ } Restrict
Lemma~\ref{lem:krein_sep_qlcms_verification} to $\Mcal^{L}$; the
boundedness condition is preserved.

\emph{Contains approximate unit:\ } By Latr\'emoli\`ere Remark~1.41,
this follows from the Leibniz property applied to the truncated
projector sequence $h_{n}$, which is M-block-diagonal hence in
$\Mcal^{L}$.
\end{proof}

\begin{definition}[Krein-pointed proper QMS]\label{def:krein_ppqms}
A \emph{Krein-pointed proper quantum metric space} (Krein PPQMS) is
a 4-tuple
\[
\mathbb{X}^{K} = (\Acal^{K}, L^{K}, \Mcal^{L}, \BWvac)
\]
where $(\Acal^{K}, L^{K}, \BWvac)$ is a Krein-pinned QLCMS,
$\Mcal^{L}$ is a Krein topography of $(\Acal^{K}, L^{K})$, $\BWvac$
restricts to a character of $\Mcal^{L}$, and there exists a Krein
L-Lipschitz $\BWvac$-pinned exhaustive sequence contained in
$\Mcal^{L}$.
\end{definition}

\begin{lemma}[The merged substrate is a Krein PPQMS]\label{lem:merged_substrate}
The 4-tuple $(\Acal^{K}, L^{K}, \Mcal^{L}, \BWvac)$ with $\Acal^{K}$
the natural Lorentzian operator system, $L^{K}$ the Krein-Leibniz
Lipschitz seminorm, $\Mcal^{L}$ the natural Krein topography of
Lemma~\ref{lem:natural_krein_topography}, and $\BWvac$ the BW vacuum,
satisfies all axioms of Definition~\ref{def:krein_ppqms}.
\end{lemma}

\begin{proof}[Proof sketch]
The (FM-K), (CB-K), (B-K) axioms hold by
Lemma~\ref{lem:krein_sep_qlcms_verification}; the continuum tightness
holds by Lemma~\ref{lem:krein_continuum_tightness}; the topography
axioms hold by Lemma~\ref{lem:natural_krein_topography};
$\BWvac|_{\Mcal^{L}}$ is a character because $\Mcal^{L}$ is Abelian
(states on Abelian C*-algebras are characters by Gelfand); the
truncated projector sequence $h_{n} = P_{\nmax(n), \Nt(n), T(n)}$ lies
in $\Mcal^{L}$ (M-block-diagonal) and is a Krein L-Lipschitz
$\BWvac$-pinned exhaustive sequence.
\end{proof}

\subsection{Krein M-tunnels with extent element}\label{sec:krein_mtunnels}

\begin{definition}[Krein M-tunnel]\label{def:krein_mtunnel}
Let $\mathbb{X}^{K}_{i} = (\Acal^{K}_{i}, L^{K}_{i}, \Mcal^{L}_{i},
\mu_{i})$ for $i = 1, 2$ be two Krein PPQMS and $M \ge 1$. A
\emph{Krein M-tunnel} from $\mathbb{X}^{K}_{1}$ to $\mathbb{X}^{K}_{2}$
is a 6-tuple
\[
\tau^{L} = (\Dcal^{L}, L_{\Dcal^{L}}^{K}, \mathfrak{M}_{\Dcal^{L}}^{L},
\pi_{1}^{K}, \pi_{2}^{K}, e^{L})
\]
where $(\Dcal^{L}, L_{\Dcal^{L}}^{K}, \mathfrak{M}_{\Dcal^{L}}^{L},
\mu_{1} \circ \pi_{1}^{K})$ is a pointed topographic separable
Krein-side QLCMS, $e^{L} \in \mathrm{dom}_{\mathrm{sa}}(L_{\Dcal^{L}}^{K})
\cap \mathfrak{M}_{\Dcal^{L}}^{L}$ is the \emph{extent element}, and
$\pi_{i}^{K} :\ \Dcal^{L} \twoheadrightarrow \Acal^{K}_{i}$ are
topographic Krein M-isometries.
\end{definition}

\textbf{The natural Krein M-tunnel.} For comparison of the
truncated Krein triple $\Tcal^{L}_{\nmax, \Nt, T}$ with the
continuum $\Tcal^{L}_{\Manifold}$:
\begin{itemize}
\item $\Dcal^{L} = \Op^{L}_{\nmax, \Nt, T} \oplus C(\Tcal^{L}_{\Manifold})$,
direct sum of truncated operator system and continuum function algebra;
\item $L_{\Dcal^{L}}^{K}((a, f)) = \max\{L^{K}_{1}(a), L^{K}_{2}(f),
\opnorm{a - \Bjoint(f)}/\varepsilon_{0}\}$, the tunnel Lipschitz
combining intrinsic seminorms with Berezin discrepancy at scale
$\varepsilon_{0} = \gammajoint$;
\item $\mathfrak{M}_{\Dcal^{L}}^{L} = \Mcal^{L}_{1} \oplus \Mcal^{L}_{2}$;
\item $\pi_{1}^{K}(a, f) = a$, $\pi_{2}^{K}(a, f) = f$;
\item $e^{L} = (h_{n}, \mathbf{1})$, the natural extent element
pairing the truncation projector with the continuum unit.
\end{itemize}

\begin{lemma}\label{lem:natural_mtunnel}
The natural Krein M-tunnel satisfies Definition~\ref{def:krein_mtunnel}.
\end{lemma}

\begin{proof}[Proof sketch]
The topography $\mathfrak{M}_{\Dcal^{L}}^{L}$ is Abelian by direct
sum of Abelian sub-algebras; $e^{L}$ is in
$\mathrm{dom}_{\mathrm{sa}}(L_{\Dcal^{L}}^{K})$ because both
$h_{n}$ and $\mathbf{1}$ are; $\pi_{1, 2}^{K}$ are topographic
M-isometries by Paper~45~\cite{paper45} Sub-Sprint~D
Berezin/projection construction; $e^{L}$ is K$^{+}$-positive because
both factors commute with $\JL$ trivially.
\end{proof}

% =====================================================================
\section{Krein-pointed propinquity
hypertopology}\label{sec:krein_hypertopology}
% =====================================================================

\subsection{Extent of a Krein M-tunnel}\label{sec:krein_extent}

\begin{definition}[Krein extent of an M-tunnel]\label{def:krein_extent}
The \emph{Krein extent} of an M-tunnel $\tau^{L} = (\Dcal^{L},
L_{\Dcal^{L}}^{K}, \mathfrak{M}_{\Dcal^{L}}^{L}, \pi_{1}^{K},
\pi_{2}^{K}, e^{L})$ is
\[
\chi^{K}(\tau^{L}) := \max\{\chi_{1}^{K}, \chi_{2}^{K}, \chi_{3}^{K}\},
\]
where $\chi_{1}^{K}$ is the reach (Hausdorff distance of pullback
quasi-state spaces), $\chi_{2}^{K}$ is the height (pin-state alignment
via Monge--Kantorovich on K$^{+}$ Wasserstein--Kantorovich), and
$\chi_{3}^{K}$ is the extent-element Lipschitz-and-normalization:
\[
\chi_{3}^{K}(\tau^{L}) = \max\{2M L^{K}(e^{L}),\
|1 - \mu_{1}\circ\pi_{1}^{K}(e^{L})|,\
|1 - \mu_{2}\circ\pi_{2}^{K}(e^{L})|\}.
\]
\end{definition}

\begin{lemma}\label{lem:extent_bound}
For the natural Krein M-tunnel from $\Tcal^{L}_{\nmax, \Nt, T}$ to
$\Tcal^{L}_{\Manifold}$,
\[
\chi^{K}(\tau^{L}) \le \gammajoint_{\nmax, \Nt, T} \to 0,
\]
bit-identical to the Paper~45~\cite{paper45} Sub-Sprint~D propinquity
bound.
\end{lemma}

\begin{proof}[Proof sketch]
$\chi_{1}^{K} = \mathrm{reach}_{B}^{K} \le \gammajoint$ by
Paper~45~\cite{paper45} Sub-Sprint~D~\S~3.1;
$\chi_{2}^{K} = $ K$^{+}$ pin-state distance, bounded by $\gammajoint$;
$\chi_{3}^{K}$ has three components, each bounded by $\gammajoint$
times a constant via the approximate-identity property of $h_{n}$
applied to the extent element.
\end{proof}

\subsection{The Krein M-local quantum metametric}\label{sec:krein_mlocal_metametric}

\begin{definition}[Krein M-local quantum metametric]\label{def:krein_mlocal}
For $M \ge 1$ and Krein PPQMS $\mathbb{X}^{K}, \mathbb{Y}^{K}$:
\[
\Dethmet_{M}(\mathbb{X}^{K}, \mathbb{Y}^{K})
:= \inf\{\chi^{K}(\tau^{L}) :\ \tau^{L} \text{ is a Krein M-tunnel from }
\mathbb{X}^{K} \text{ to } \mathbb{Y}^{K}\}.
\]
\end{definition}

\begin{theorem}[4-point relaxed forward triangle inequality]\label{thm:relaxed_triangle}
For any $M \ge 1$ and Krein PPQMS $\mathbb{X}^{K}, \mathbb{Y}^{K},
\mathbb{Z}^{K}$:
\[
\Dethmet_{M}(\mathbb{X}^{K}, \mathbb{Y}^{K}) \le
(1 + \Dethmet_{M}(\mathbb{X}^{K}, \mathbb{Z}^{K}))^{2}\,
\Dethmet_{M}(\mathbb{Z}^{K}, \mathbb{Y}^{K}) +
(1 + \Dethmet_{M}(\mathbb{Z}^{K}, \mathbb{Y}^{K}))^{2}\,
\Dethmet_{M}(\mathbb{X}^{K}, \mathbb{Z}^{K}).
\]
\emph{(Structural/categorical: $\Dethmet_{M}$ is built from the
degenerate Krein Lipschitz seminorm of Paper~45, so this is a property
of the construction; its metric-level reading is descoped pending
substrate repair --- see the Status note.)}
\end{theorem}

\begin{proof}[Proof sketch]
Direct transport of Latr\'emoli\`ere 2512.03573 Thm~2.24 via Krein-side
tunnel composition. The natural-substrate Krein structure is closed
under tunnel composition (direct sums of operator systems give operator
systems, Berezin composition preserves K$^{+}$-positivity, the
extent-element construction generalizes to triple composition).
\end{proof}

\subsection{Coincidence}\label{sec:coincidence}

\begin{theorem}[Coincidence]\label{thm:coincidence}
For Krein PPQMS $\mathbb{X}^{K}, \mathbb{Y}^{K}$, there exists a
full topographic Krein M-isometry $\pi^{K} : \mathbb{X}^{K} \to
\mathbb{Y}^{K}$ with $\mu_{\mathbb{Y}^{K}} \circ \pi^{K} =
\mu_{\mathbb{X}^{K}}$ if and only if
$\Dethmet_{M}(\mathbb{X}^{K}, \mathbb{Y}^{K}) = 0$.
\emph{(Descoped at the metric level: with the degenerate Krein seminorm
of Paper~45 the ``only if'' direction $\Dethmet_{M}=0\Rightarrow$
isometry is conditional on substrate repair; as it stands the statement
is categorical/structural --- see the Status note.)}
\end{theorem}

\begin{proof}[Proof sketch]
The "if" direction is contained in Theorem~\ref{thm:relaxed_triangle}.
The "only if" direction transports Latr\'emoli\`ere 2512.03573
Thm~3.6 via target-set machinery:\ given a sequence of tunnels
$\tau_{n}^{K}$ with $\chi^{K}(\tau_{n}^{K}) \to 0$, the target sets
$\mathfrak{t}_{\tau^{K}}^{K}(a | l) :=
\{\pi_{2}^{K}(d) : \pi_{1}^{K}(d) = a,\ \norm{d}_{L_{\Dcal^{L}}, M}
\le l\}$ are well-defined on the Krein side because of the
K$^{+}$-positivity of natural-substrate generators, and the four
almost-properties (Latr\'emoli\`ere Lemmas~3.2--3.5) transport
mechanically via Krein-Leibniz and Krein-approximate-unit. The
diagonal extraction with Banach--Alaoglu on the K$^{+}$-restricted
state space produces a limit *-isomorphism preserving topography and
BW pin state.
\end{proof}

\subsection{The Krein quantum metametric and
hypertopology}\label{sec:krein_hypertopology_metametric}

\begin{definition}[Krein quantum metametric]\label{def:krein_metametric}
For Krein PPQMS $\mathbb{X}^{K}, \mathbb{Y}^{K}$:
\[
\BigDeth(\mathbb{X}^{K}, \mathbb{Y}^{K}) := \max\bigl\{\,
\inf\{\varepsilon > 0 : \sup_{r \in [1, 1/\varepsilon]}
\Dethmet_{r}(\mathbb{X}^{K}, \mathbb{Y}^{K}) < \varepsilon\},\
\bigl|e^{-\mathrm{qdiam}(\mathbb{X}^{K})} -
e^{-\mathrm{qdiam}(\mathbb{Y}^{K})}\bigr|\,\bigr\}.
\]
\end{definition}

\begin{theorem}[Inframetric structure]\label{thm:inframetric}
$\BigDeth$ is a K$^{+}$-positive inframetric:\ symmetric, satisfying
the 4-point relaxed triangle, and $\BigDeth(\mathbb{X}^{K},
\mathbb{Y}^{K}) = 0$ iff a full topographic Krein M-isometry exists
between $\mathbb{X}^{K}$ and $\mathbb{Y}^{K}$.
\emph{(Structural/categorical: built on the degenerate Krein seminorm
of Paper~45; the non-degeneracy (metric-level) reading is descoped
pending substrate repair --- see the Status note.)}
\end{theorem}

\begin{proof}[Proof sketch]
Symmetry holds because Krein M-tunnels are symmetric in their roles;
the 4-point relaxed triangle follows from
Theorem~\ref{thm:relaxed_triangle}; the coincidence statement is
Theorem~\ref{thm:coincidence}.
\end{proof}

\begin{definition}[Krein hypertopology]\label{def:krein_hypertopology}
The \emph{Krein hypertopology} on the class
$p\mathcal{P}^{K}$ of Krein PPQMS is the topology with Kuratowski
closure operator
\[
\mathrm{cl}^{K}(\mathcal{A}) := \{\mathbb{Y}^{K} \in p\mathcal{P}^{K} :
\BigDeth(\mathbb{Y}^{K}, \mathcal{A}) = 0\}
\]
where $\BigDeth(\mathbb{Y}^{K}, \mathcal{A}) :=
\inf\{\BigDeth(\mathbb{Y}^{K}, \mathbb{X}^{K}) : \mathbb{X}^{K}
\in \mathcal{A}\}$.
\end{definition}

\subsection{Compact-case agreement at $\Nt = 1$}\label{sec:compact_case}

\begin{theorem}[Compact-case agreement]\label{thm:compact_agreement}
When restricted to $\Nt = 1$, the Krein operator-system substrate and
the SU(2) rate factor reduce bit-exactly to the Paper~38~\cite{paper38}
SU(2) state-space Gromov--Hausdorff construction. The \emph{joint} L1
rate carries an additive U(1) temporal offset $T/4$ (the Riemannian-limit
temporal factor), so the bit-exact reduction is on the spatial substrate
and its SU(2) rate, not on the joint rate.
\end{theorem}

\begin{proof}[Proof sketch]
At $\Nt = 1$:\ $\Krein_{\nmax, 1, T} = \HGV^{\nmax}$ (spatial Krein
space alone); $\JL = J_{\spat}$; $\Kplus = \HGV^{\nmax, +}$;
$\DL|_{\Nt=1} = i\DGV$ (temporal direction contributes nothing);
$\Op^{L}|_{\Nt=1}$ is the chirality-doubled spinor lift of
Paper~38~\cite{paper38}'s scalar substrate. By Paper~45~\cite{paper45}
Sub-Sprint~D~\S~5 Riemannian-limit recovery, the SU(2) rate factor
$\gamma_{\mathrm{SU(2)}}(\nmax) = \Lambda^{\mathrm{Paper~38}}_{\nmax}$
bit-exact; the joint rate is $\gamma_{\mathrm{L1}}(\nmax, 1, T) =
\gamma_{\mathrm{SU(2)}}(\nmax) + T/4$, the $T/4$ being the U(1)
temporal-factor offset (uniform-circle expected $|\theta|$). The
operator-system / spatial-substrate reduction is therefore bit-exact;
the joint rate matches in its SU(2) factor up to that $T/4$ offset.
\end{proof}

\begin{remark}[The compact-case agreement object]\label{rem:compact_object}
The compact-case object is the chirality-doubled spinor lift of
Paper~38's Riemannian SU(2) Lipschitz spectral triple, pinned at the
BW vacuum on the spatial wedge. This explicit identification ensures
the Krein-pointed hypertopology is a strict extension of
Paper~38's state-space Gromov--Hausdorff construction, not a parallel one.
\end{remark}

\subsection{Numerical panel}\label{sec:numerical_panel}

The Krein-side propinquity bound inherits the bit-exact numerical
panel from Paper~45~\cite{paper45}:

\begin{center}
\begin{tabular}{ccc}
\toprule
$(\nmax, \Nt)$ & $\Dethmet_{M}$ (qualitative) & $\Lprop$ (Paper~45) \\
\midrule
$(2, 3)$ & $\le 2.0746$ & $= 2.0746$ \\
$(3, 5)$ & $\le 1.6101$ & $= 1.6101$ \\
$(4, 7)$ & $\le 1.3223$ & $= 1.3223$ \\
\bottomrule
\end{tabular}
\end{center}

Monotone decreasing, with ratio $\Dethmet(4, 7)/\Dethmet(2, 3) = 0.6374$
matching Paper~38~\cite{paper38} single-factor bit-identical.

% =====================================================================
\section{The Mondino--S\"amann bridge:\ Wick-rotation
functor}\label{sec:bridge}
% =====================================================================

\subsection{The F2 mismatch and three candidate
resolutions}\label{sec:f2_resolution}

The structural mismatch between the Krein metametric $\BigDeth$ (4-point
relaxed forward, codomain $[0, \infty]$) and the Mondino--S\"amann
time separation $\ell$ (3-point reverse, codomain $\{-\infty\}\cup
[0, \infty]$) is the \textbf{F2 mismatch}. We test three candidate
resolutions.

\subsubsection{R1:\ Negate time}\label{sec:r1_test}

\textbf{Candidate.} Define $\ell^{L}_{R1}(\hat{\omega}, \hat{\omega}')
:= -\BigDeth(\mathbb{X}^{K}_{\hat{\omega}}, \mathbb{X}^{K}_{\hat{\omega}'})$.

\textbf{Verdict:\ NEGATIVE.} R1 fails for two independent reasons:
(i) Codomain mismatch:\ $\BigDeth \in [0, \infty]$ gives
$\ell^{L}_{R1} \in [-\infty, 0]$, while Mondino--S\"amann requires
$\{-\infty\}\cup [0, \infty]$. The MS reverse triangle encodes the
Lorentzian convention that "longer causal curves accumulate more
proper time" and the negation gives the opposite convention.
(ii) Triangle structural mismatch:\ negating both sides of the
4-point relaxed forward inequality does NOT collapse to the 3-point
reverse triangle, and the coefficient factors $(1 + \BigDeth)^{2}$
have no MS-side analog.

\subsubsection{R2:\ Functorial composite of Wick-rotation +
modular-flow}\label{sec:r2_test}

\textbf{Candidate.} A bridge functor $F_{a} \circ F_{c}$ composed of:\
(a) Wick rotation $F_{a}$ giving the right codomain shape (3-point
reverse triangle on boost orbits via analytic continuation), and
(c) modular flow $F_{c}$ giving the right basepoint (BW vacuum) and
the right cover scales ($\beta_{k} = 2\pi/\kappag$).

\textbf{Verdict:\ PARTIAL.} (a) gives the right structural shape
but does not determine basepoint or cover scales. (c) gives basepoint
and scales but inherits the F2 triangle-direction mismatch (modular
flow is additive, not super-additive). The composite $F_{a} \circ F_{c}$
has all needed ingredients but the F2 mismatch is not resolved
\emph{by R2 alone}; the resolution requires R3.

\subsubsection{R3:\ Connes--Rovelli thermal time as unifying
framework}\label{sec:r3_test}

\textbf{Candidate.} Identify the Krein metametric $\BigDeth$ as
measuring \emph{thermal time} of the wedge KMS state and the
Mondino--S\"amann $\ell$ as measuring \emph{geometric proper time} of
the same wedge KMS state. The Wick rotation between them, via
equation~\eqref{eq:connes_rovelli_identification}, effects the
analytic continuation $t_{\mathrm{thermal}} = -it_{\mathrm{geometric}}/\kappag$
mapping additive thermal-time structure to super-additive
geometric-time structure.

\textbf{Verdict:\ POSITIVE.} R3 supplies the unifying interpretation
that resolves F2. Modular flow $\sigma_{t_{1}}^{\omega} \circ
\sigma_{t_{2}}^{\omega} = \sigma_{t_{1} + t_{2}}^{\omega}$ is
\emph{additive} in thermal time (forward-triangle), while geometric
proper time on a Lorentzian manifold is \emph{super-additive} on
causal curves (the "twin paradox"). The two triangle directions are
not "mismatched":\ they encode the same wedge KMS state under two
distinct projection mechanisms (operator-algebraic modular flow vs.\
synthetic causal geometry), connected by the Wick rotation.

\subsection{Definition of the Wick-rotation functor}\label{sec:wick_functor_def}

\begin{definition}[The category $\KreinMM$]\label{def:cat_kreinmm}
Let $\KreinMM$ be the category whose:
\begin{itemize}
\item Objects are Krein PPQMS $\mathbb{X}^{K} = (\Acal^{K}, L^{K},
\Mcal^{L}, \BWvac)$ per Definition~\ref{def:krein_ppqms};
\item Morphisms are topographic Krein M-isometries
$\pi^{K} : \mathbb{X}^{K}_{1} \to \mathbb{X}^{K}_{2}$, i.e.\ proper
*-epimorphisms preserving topography and pin state.
\end{itemize}
\end{definition}

\begin{definition}[The category $\LorPLG$]\label{def:cat_lorplg}
Let $\LorPLG$ be the category whose:
\begin{itemize}
\item Objects are covered LPLS $(X, \ell, o, \Ucal)$ per
Definition~\ref{def:covered_lpls};
\item Morphisms are $\ell$-preserving maps preserving basepoint $o$
and cover $\Ucal$.
\end{itemize}
\end{definition}

\begin{definition}[The Wick-rotation functor]\label{def:wick_functor}
The \emph{Wick-rotation functor}
$\Wfun : \KreinMM \to \LorPLG$ is defined on objects by
\[
\Wfun(\Acal^{K}, L^{K}, \Mcal^{L}, \BWvac)
:= (\hat{\Mcal}^{L}, \ell^{L}, \hat{\omega}_{W}^{L}, \hat{\Ucal}^{L}),
\]
where:
\begin{enumerate}
\item $\hat{\Mcal}^{L} := \mathrm{Spec}(\Mcal^{L})$ is the Gelfand
spectrum of the M-diagonal Abelian topography;
\item $\hat{\omega}_{W}^{L}$ is the character of $\Mcal^{L}$ induced
by the BW vacuum;
\item $\hat{\Ucal}^{L} = (\hat{U}_{k})_{k \in \N}$ with
$\hat{U}_{k} := \mathrm{Spec}(\Mcal^{L}|_{(\nmax(k), \Nt(k), T(k))})$
along admissible scaling per Paper~47~\cite{paper47};
\item The time separation $\ell^{L}$ on $\hat{\Mcal}^{L}$ is
\[
\ell^{L}(\hat{\omega}, \hat{\omega}') := \begin{cases}
\kappag \cdot \tmod^{\BWvac}(\hat{\omega}, \hat{\omega}') &
\text{if } \hat{\omega} \preceq \hat{\omega}' \text{ along modular flow,} \\
-\infty & \text{otherwise.}
\end{cases}
\]
\end{enumerate}
On morphisms, $\Wfun(\pi^{K}) := \hat{\pi}^{K}$ is the dual Gelfand
spectrum map (contravariant by Gelfand duality:\
$\KreinMM \to \LorPLG$ on objects sends $\mathbb{X}^{K}$ forward but
on morphisms the arrow direction reverses).
\end{definition}

\subsection{The Krein--Mondino--S\"amann Bridge
Theorem}\label{sec:bridge_theorem}

\begin{theorem}[Krein--Mondino--S\"amann
Bridge]\label{thm:bridge}
The Wick-rotation functor $\Wfun$ of Definition~\ref{def:wick_functor}
is well-defined, with the following four properties:
\begin{itemize}
\item[\textbf{(B1)}] \emph{Structural correspondence.} $\Wfun$ sends
the Krein PPQMS substrate to the Mondino--S\"amann covered LPLS
structure row-by-row as in Table~\ref{tab:bridge_correspondence}.
\item[\textbf{(B2)}] \emph{Reverse triangle inequality.} For every
$\mathbb{X}^{K}$ and every $\hat{\omega}_{x}, \hat{\omega}_{y},
\hat{\omega}_{z} \in \Wfun(\mathbb{X}^{K})$,
\[
\ell^{L}(\hat{\omega}_{x}, \hat{\omega}_{y}) +
\ell^{L}(\hat{\omega}_{y}, \hat{\omega}_{z}) \le
\ell^{L}(\hat{\omega}_{x}, \hat{\omega}_{z}).
\]
\item[\textbf{(B3)}] \emph{Pre-compactness inheritance.} The
Mondino--S\"amann pre-compactness conditions (Thm~6.2 of MS) hold on
$\Wfun(\mathbb{X}^{K}_{n})$ for any sequence of truncated Krein PPQMS
$(\mathbb{X}^{K}_{n})$ at admissible-scaling cutoffs.
\item[\textbf{(B4)}] \emph{Convergence transport} (\textbf{descoped}). If
$\BigDeth(\mathbb{X}^{K}_{n}, \mathbb{X}^{K}_{\infty}) \to 0$, then
$\Wfun(\mathbb{X}^{K}_{n}) \to \Wfun(\mathbb{X}^{K}_{\infty})$ in the
Mondino--S\"amann pLGH sense. \emph{(The premise $\BigDeth\to0$ is built
from the degenerate Krein metric of Paper~45, so this transport is
conditional pending substrate repair, not an established metric
convergence --- see Theorem~\ref{thm:convergence_transport}.)}
\end{itemize}
\end{theorem}

\subsection{Structural correspondence (B1)}\label{sec:b1_correspondence}

The structural correspondence holds row-by-row per
Table~\ref{tab:bridge_correspondence}:

\begin{table}[h]
\centering
\caption{The structural correspondence (B1) between Mondino--S\"amann
covered LPLS and Krein PPQMS structures, exhibited by the
Wick-rotation functor $\Wfun$.}
\label{tab:bridge_correspondence}
\begin{tabular}{rll}
\toprule
\# & Mondino--S\"amann (synthetic) & Krein-algebraic-side analog \\
\midrule
1 & Underlying set $X$ & $\hat{\Mcal}^{L} = \mathrm{Spec}(\Mcal^{L})$ \\
2 & Topology finer than chronological & weak-$*$ topology on $\hat{\Mcal}^{L}$ \\
3 & Time separation $\ell$ & Wick-rotated $\ell^{L} = \kappag\cdot\tmod$ \\
4 & Reverse triangle (Eq.~\eqref{eq:reverse_triangle}) & on orbit:\ equality; off orbit:\ trivial $-\infty$ \\
5 & Timelike $\ll$ & modular precedence with $t > 0$ \\
6 & Causal $\le$ & modular precedence with $t \ge 0$ \\
7 & Causal diamond & $\{\hat{\omega}_{x}\circ\sigma_{s} : 0 \le s \le t\}$ \\
8 & Basepoint event $o$ & BW vacuum $\hat{\omega}_{W}^{L}$ \\
9 & Cover $\Ucal$ & truncated topography spectra $(\hat{U}_{k})$ \\
10 & Cover scales $\beta_{k}$ & BW canonical period $\beta_{k} = 2\pi/\kappag$ \\
11 & Slab bound & uniform $2\pi$ via modular periodicity \\
12 & ε-net & modular causal diamond ε-net at scale $\gammajoint(k)$ \\
13 & LGH convergence & truncation-projector correspondence \\
14 & pLGH convergence & per-cover LGH convergence \\
15 & Pre-compactness & cardinality bound from propagation $= 2$ \\
\bottomrule
\end{tabular}
\end{table}

\subsection{Reverse triangle inequality (B2):\ Decomposition O on the
M-diagonal topography}\label{sec:b2_decomposition_o}

The substantive content of (B2) is the resolution of F2 via R3 plus
the structural simplification that off-orbit cases are trivial at the
K$^{+}$-weak-form bridge.

\begin{lemma}[Modular flow acts trivially on the
topography]\label{lem:trivial_modular_action_on_topography}
For every topographic generator $a = M^{\spat}_{N, L, 0} \otimes I_{\Nt}
\in \Mcal^{L}$,
\[
\sigma_{t}^{\BWvac}(a) = a \quad \forall t \in \R.
\]
Consequently, the dual action on $\hat{\Mcal}^{L}$ is trivial:\
$\hat{\sigma}_{t}^{\BWvac}(\hat{\omega}) = \hat{\omega}$ for all
$\hat{\omega}$.
\end{lemma}

\begin{proof}
$[\Kalpha, M^{\spat}_{N, L, 0}] = 0$ because $\Kalpha$ is M-diagonal
and the topographic generator has $M = 0$, both diagonal in the same
$(N, L, M)$-block decomposition. Hence $\sigma_{t}^{\BWvac}(a) =
e^{-it\Kalpha} a e^{+it\Kalpha} = a$, and the dual Gelfand action is
trivial.
\end{proof}

\begin{lemma}[Decomposition O for triples of
characters]\label{lem:decomposition_o}
For $\hat{\omega}_{x}, \hat{\omega}_{y}, \hat{\omega}_{z} \in
\hat{\Mcal}^{L}$ at finite cutoff, indexed by topographic eigenvalue
labels $(N_{x}, L_{x}), (N_{y}, L_{y}), (N_{z}, L_{z})$ with $M = 0$,
the off-orbit case decomposes into four sub-cases:
\begin{itemize}
\item[(i)] All three on the same modular-flow orbit:\
$(N_{x}, L_{x}) = (N_{y}, L_{y}) = (N_{z}, L_{z})$. On-orbit case;
\item[(ii)] $\hat{\omega}_{x}, \hat{\omega}_{z}$ on the same orbit,
$\hat{\omega}_{y}$ on a different orbit. Then
$\ell^{L}(\hat{\omega}_{x}, \hat{\omega}_{y}) = -\infty$ (or
$\ell^{L}(\hat{\omega}_{y}, \hat{\omega}_{z}) = -\infty$);
\item[(iii)] $\hat{\omega}_{x}, \hat{\omega}_{y}$ on one orbit,
$\hat{\omega}_{y}, \hat{\omega}_{z}$ on a different orbit,
$\hat{\omega}_{x}, \hat{\omega}_{z}$ on yet a third. \textbf{Empty}
(this would require $(N_{x}, L_{x}) = (N_{y}, L_{y}) \ne
(N_{z}, L_{z}) = (N_{x}, L_{x})$, a contradiction);
\item[(iv)] All three on different orbits. At least one
$\ell^{L} = -\infty$.
\end{itemize}
\end{lemma}

\begin{proof}[Proof sketch]
The orbit-preservation of modular flow (Lemma~\ref{lem:trivial_modular_action_on_topography})
means modular flow on $\hat{\Mcal}^{L}$ preserves the $(N, L)$ labels.
Hence two characters are causally related ($\hat{\omega} \preceq
\hat{\omega}'$ along modular flow with finite $\tmod$) only if they
share the same $(N, L)$ label. Cases (i)--(iv) are then a complete
enumeration of label-tuple configurations, and (iii) is structurally
empty.
\end{proof}

\begin{proof}[Proof of Theorem~\ref{thm:bridge} (B2)]
By Lemma~\ref{lem:decomposition_o}, we analyze (B2) case by case.

\emph{Case (i), on-orbit (equality).} If $\hat{\omega}_{x} \preceq
\hat{\omega}_{y} \preceq \hat{\omega}_{z}$ along modular flow with
$\hat{\omega}_{y} = \hat{\omega}_{x}\circ\sigma_{t_{1}}^{\BWvac}$ and
$\hat{\omega}_{z} = \hat{\omega}_{y}\circ\sigma_{t_{2}}^{\BWvac}$,
then $\hat{\omega}_{z} = \hat{\omega}_{x}\circ\sigma_{t_{1} +
t_{2}}^{\BWvac}$ by additivity of modular flow. Hence
$\ell^{L}(\hat{\omega}_{x}, \hat{\omega}_{y}) + \ell^{L}(\hat{\omega}_{y},
\hat{\omega}_{z}) = \kappag(t_{1} + t_{2}) = \ell^{L}(\hat{\omega}_{x},
\hat{\omega}_{z})$.

\emph{Cases (ii), (iv), trivial $-\infty$ inheritance.} At least one
of $\ell^{L}(\hat{\omega}_{x}, \hat{\omega}_{y})$ or
$\ell^{L}(\hat{\omega}_{y}, \hat{\omega}_{z})$ is $-\infty$, so the
LHS is $-\infty$ by MS convention, and $-\infty \le
\ell^{L}(\hat{\omega}_{x}, \hat{\omega}_{z})$ trivially.

\emph{Case (iii), empty.} No contribution.

Hence (B2) holds in all cases.
\end{proof}

\begin{remark}[The super-additivity that DOES exist lives in
$\Op^{L}$, not in $\Mcal^{L}$]\label{rem:super_add_lives_elsewhere}
The substantive super-additivity property (off-axis Lorentz-boost
composition accumulating more proper time than direct composition)
exists on the wedge but lives in the \emph{full} Krein operator
system $\Op^{L}$, not in the M-diagonal topography $\Mcal^{L}$.
The Paper~46~\cite{paper46} Appendix~B enlarged substrate (chirality-
flipping generators $M^{\mathrm{flip}}$ with $\{\JL, M^{\mathrm{flip}}\}
= 0$) supports such off-axis composition; on this enlarged substrate,
the bridge $\Wfun$ would extend to a richer LPLS structure with strict
super-additivity off-orbit. This is the open question Q1'
(Section~\ref{sec:open_q1prime}).
\end{remark}

\subsection{Pre-compactness inheritance (B3)}\label{sec:b3_precompactness}

\begin{proof}[Proof of Theorem~\ref{thm:bridge} (B3)]
The Mondino--S\"amann pre-compactness Thm~6.2 has three conditions:
(i) uniform slab bounds; (ii) uniform ε-net cardinality bounds;
(iii) cumulative ε-net nesting.

\emph{(i) Uniform slab bound.} The timelike diameter of $\hat{U}_{k}$
equals $\sup \tmod^{\BWvac}$ on the wedge. Under BW canonical
$\kappag = 1$, modular flow is $2\pi$-periodic (Paper~42~\cite{paper42}
\S~5 integer spectrum of $\Kalpha$), so $\tmod \le 2\pi$ uniformly:\
$T_{k} = 2\pi$ for all $k$.

\emph{(ii) Uniform ε-net cardinality bound.} At finite cutoff
$(\nmax(k), \Nt(k), T(k))$, $|\hat{U}_{k}| = O(\nmax(k)^{4} \cdot \Nt(k))$
(dimension of the topography), and the ε-net cardinality is bounded
above by $|\hat{U}_{k}|$ itself and by the standard discrete-to-continuous
refinement bound $N(k, \varepsilon) = \min(|\hat{U}_{k}|,
\mathrm{Vol}(\hat{U}_{k})/\varepsilon^{d_{k}})$. Both bounds are
uniform via Paper~44~\cite{paper44} propagation number $= 2$.

\emph{(iii) Cumulative ε-net nesting.} By the admissible-scaling
sequence $(\nmax(k), \Nt(k), T(k)) \to (\infty, \infty, T_{\infty})$,
the cover $\hat{U}_{k} \subseteq \hat{U}_{k+1}$ is nested by
construction, and ε-nets of $\hat{U}_{k}$ extend to ε-nets of
$\hat{U}_{k+1}$ by adding finitely many points.

By Mondino--S\"amann Thm~6.2, the bridge image inherits
pre-compactness.
\end{proof}

\subsection{Convergence transport (B4) via Plancherel-nested
Berezin}\label{sec:b4_convergence}

\begin{theorem}[Convergence transport --- \emph{descoped}]\label{thm:convergence_transport}
If $\BigDeth(\mathbb{X}^{K}_{n}, \mathbb{X}^{K}_{\infty}) \to 0$ in
the Krein hypertopology of Definition~\ref{def:krein_hypertopology},
then $\Wfun(\mathbb{X}^{K}_{n}) \to \Wfun(\mathbb{X}^{K}_{\infty})$
in the Mondino--S\"amann pLGH topology of
Definition~\ref{def:plgh}, with the strong (timelike forward density)
sense. \emph{(Descoped:} $\BigDeth$ is built from the degenerate Krein
Lipschitz seminorm $L^{K}$ of Paper~45, so the premise $\BigDeth\to0$ is
not a non-degenerate metric convergence; this transport is conditional
pending substrate repair, not an established metric-level result.\emph{)}
\end{theorem}

\begin{proof}[Proof sketch]
By the Krein hypertopology coincidence (Theorem~\ref{thm:coincidence})
plus the bit-exact panel inheritance from Paper~45~\cite{paper45},
the Krein-side propinquity convergence implies that for each cover
level $k$, the truncated topography spectra $\hat{U}_{k, n} \to
\hat{U}_{k, \infty}$ in the LGH sense (per
Definition~\ref{def:lgh_convergence}). The Plancherel-nested
Berezin/projection pair from Paper~45~\cite{paper45} Sub-Sprint~C/D
supplies a Mondino--S\"amann correspondence with vanishing
distortion. Per-axiom verification:

\emph{Cardinality matching (MS Def~3.6(i)).} At fixed admissible
cutoff, the correspondence pair has $|S_{n}| = |S_{\infty}|$ on the
M-diagonal topography (the Berezin map preserves the $(N, L)$ shell
structure).

\emph{Correspondence with vanishing distortion (MS Def~3.6(ii)).}
The Berezin/projection pair has reach $\le \gammajoint$ + height $\le
\gammajoint \cdot C_{3}^{\mathrm{joint}}$, vanishing as $n \to \infty$.

\emph{Extension property (MS Def~3.6(iii)).} The Plancherel-nested
structure of Berezin maps IS the MS extension property:\ Berezin maps
at finer cutoff extend Berezin maps at coarser cutoff because the
Plancherel weights nest by construction (Paper~45~\cite{paper45}
Sub-Sprint~C Lemma~5.1).

\emph{Forward density (MS Def~3.6(iv)).} The modular-flow orbit
discretization in the chronological topology provides the timelike
forward density:\ for any $\hat{\omega}' \succ \hat{\omega}$ on a
modular orbit, the truncation projects $\hat{\omega}'$ to a sequence
$\hat{\omega}'_{n}$ converging to $\hat{\omega}'$ in the chronological
topology with $\hat{\omega} \prec \hat{\omega}'_{n}$ for all
sufficiently large $n$.

Per-cover LGH convergence at each $k$ implies pLGH convergence per
Definition~\ref{def:plgh}.
\end{proof}

\subsection{Riemannian-limit recovery cross-check}\label{sec:riemannian_limit_cross}

\begin{corollary}\label{cor:riemannian_limit_cross}
At $\Nt = 1$, the Krein-side propinquity bound $\Lprop_{\nmax, 1, T} =
\Lambda^{\mathrm{Paper~38}}_{\nmax}$ bit-exact (Paper~45 Sub-Sprint~D
Lemma~5.1), so the bridge image $\Wfun(\mathbb{X}^{K}|_{\Nt = 1})$ is
the synthetic-side counterpart of Paper~38's SU(2) state-space
Gromov--Hausdorff convergence:\ a trivial LPLS with no Lorentzian content (modular flow
on the Riemannian limit collapses to identity, and $\ell^{L}|_{\Nt=1}
= 0$ along the Gelfand spectrum).
\end{corollary}

\subsection{Honest scope of the bridge theorem}\label{sec:bridge_scope}

The bridge functor $\Wfun$ is constructed at the K$^{+}$-weak-form
level on the natural chirality-doubled scalar-multiplier substrate of
Papers~44/45. The strong-form extension of $\Wfun$ to the
Paper~46~\cite{paper46} Appendix~B enlarged substrate
(chirality-flipping generators) is \textbf{not} constructed in this
paper; it is the open question Q1' (Section~\ref{sec:open_q1prime}).

The bridge is \textbf{functorial, not isometric}:\ $\BigDeth$ and
$\ell^{L}$ measure different physical quantities (thermal vs.\
geometric time) of the same wedge KMS state. The two triangle
directions are not "mismatched" — they encode different projections
of the same physical object, connected by the Wick rotation under R3
(Connes--Rovelli thermal time).

% =====================================================================
\section{Application to the GeoVac hemispheric
wedge}\label{sec:wedge_application}
% =====================================================================

\subsection{The GeoVac wedge as a Krein PPQMS}\label{sec:wedge_as_krein_ppqms}

\begin{lemma}[The wedge as Krein PPQMS]\label{lem:wedge_as_krein_ppqms}
The Paper~43~\cite{paper43} hemispheric wedge $W_{L} = P_{W}^{\spat}
\otimes P_{t \ge 0}$ at every finite cutoff $(\nmax, \Nt, T)$ defines
a Krein PPQMS
\[
\mathbb{X}^{K}_{\mathrm{wedge}} := (\Acal^{L}|_{W_{L}},
L^{K}|_{W_{L}}, \Mcal^{L}|_{W_{L}}, \BWvac),
\]
with all five Krein PPQMS axioms of
Definition~\ref{def:krein_ppqms} verified by direct construction.
\end{lemma}

\begin{proof}[Proof sketch]
(A1) $\Acal^{L}|_{W_{L}}$ is the wedge-restricted Lorentzian operator
system (Paper~44~\cite{paper44} \S~2). (A2) $L^{K}|_{W_{L}}$ is the
K$^{+}$-restricted Krein-Leibniz Lipschitz seminorm of
Definition~\ref{def:krein_leibniz}. (A3) $\Mcal^{L}|_{W_{L}}$ is the
wedge-restricted M-diagonal Abelian topography of
Lemma~\ref{lem:natural_krein_topography}. (A4) $\BWvac =
e^{-\Kalpha}/Z$ is the wedge KMS state (Paper~43~\cite{paper43} \S~4.2),
diagonal in the M-eigenbasis hence a character of $\Mcal^{L}|_{W_{L}}$.
(A5) The Krein M-tunnel structure with $\chi^{K} \to 0$ is the wedge
restriction of Lemma~\ref{lem:extent_bound}.
\end{proof}

\subsection{The seven newly accessible theorems}\label{sec:seven_theorems}

We apply the bridge functor $\Wfun$ to the wedge PPQMS:
\[
\Wfun(\mathbb{X}^{K}_{\mathrm{wedge}}) =
(\hat{\Mcal}^{L}|_{W_{L}}, \ell^{L}_{\mathrm{wedge}},
\hat{\omega}_{W}^{L}, \hat{\Ucal}^{L}_{\mathrm{wedge}}),
\]
where at panel cells the cardinality is $|\hat{\Mcal}^{L}|_{W_{L}}|
= \dim W_{L} = 16, 60, 160$ at $(\nmax, \Nt) \in \{(2, 3), (3, 5),
(4, 7)\}$ respectively.

\subsubsection{T1:\ The wedge as Mondino--S\"amann covered
LPLS}\label{sec:t1_wedge_as_lpls}

\begin{theorem}[T1:\ GeoVac wedge as Mondino--S\"amann pointed
pre-length space]\label{thm:wedge_as_lpls}
For every finite cutoff $(\nmax, \Nt, T)$ with $T = 2\pi$ canonical,
the GeoVac hemispheric wedge $W_{L}$ defines a Mondino--S\"amann
covered Lorentzian pre-length space
\[
(\hat{\Mcal}^{L}|_{W_{L}}, \ell^{L}_{\mathrm{wedge}},
\hat{\omega}_{W}^{L}, \hat{\Ucal}^{L}_{\mathrm{wedge}})
\in \LorPLG,
\]
with uniform timelike diameter $2\pi$ across all cover levels.
\end{theorem}

\begin{proof}[Proof sketch]
Direct application of Theorem~\ref{thm:bridge} (B1) plus
Lemma~\ref{lem:wedge_as_krein_ppqms}. The uniform timelike diameter
follows from Paper~43~\cite{paper43} Thm~bw\_alpha\_lorentz:\ modular
flow is $2\pi$-periodic on the wedge.
\end{proof}

\subsubsection{T2:\ Synthetic compactness of the wedge
family}\label{sec:t2_synthetic_compactness}

\begin{theorem}[T2:\ GeoVac wedge synthetic
compactness]\label{thm:synthetic_compactness}
Any sequence $(\mathbb{X}^{K}_{\mathrm{wedge}}(k_{n}))_{n \in \N}$
of GeoVac wedge Krein PPQMS at admissible-scaling cutoffs
$(\nmax(k_{n}), \Nt(k_{n}), T(k_{n})) \to (\infty, \infty,
T_{\infty})$ has a subsequence whose bridge images
$\Wfun(\mathbb{X}^{K}_{\mathrm{wedge}}(k_{n}))$ converge in the
Mondino--S\"amann pLGH sense to a covered LPLS limit. \emph{(The
pre-compactness leg survives via the propagation-number cardinality
bound; the \textbf{convergence} of the subsequence to its limit inherits
the metric-level descope of Theorem~\ref{thm:convergence_transport} --- it
is conditional pending substrate repair.)}
\end{theorem}

\begin{proof}[Proof sketch]
By Theorem~\ref{thm:bridge} (B3), the three pre-compactness conditions
of MS Thm~6.2 hold for the wedge sequence. The MS Thm~6.2 yields
the pLGH-convergent subsequence.
\end{proof}

\subsubsection{T3:\ Synthetic bit-exact panel
inheritance}\label{sec:t3_bit_exact_panel}

\emph{[Descoped (see the Status note): the inherited values are
evaluations of the closed-form rate formula, not measurements of a
metric; the panel transports formula values, not metric data.]}

\begin{theorem}[T3:\ Synthetic pLGH-convergence rate panel for the
GeoVac wedge]\label{thm:bit_exact_panel}
The Mondino--S\"amann pLGH-convergence rate of the GeoVac wedge
sequence at the Paper~45~\cite{paper45} panel cells satisfies:

\begin{center}
\begin{tabular}{cc}
\toprule
$(\nmax, \Nt)$ & $d_{\mathrm{pLGH}}^{\mathrm{wedge}}$ (Krein-side $\Lprop$) \\
\midrule
$(2, 3)$ & $\le 2.0746$ \\
$(3, 5)$ & $\le 1.6101$ \\
$(4, 7)$ & $\le 1.3223$ \\
\bottomrule
\end{tabular}
\end{center}

bit-identical to the Paper~45 Krein-side $\Lprop$ values. These are the
descoped Paper~45 \textbf{rate-formula} evaluations (not metric-convergence
measurements), transported conditionally through the bridge:\ the
Lorentzian quantum-metric convergence is descoped, so this is a
\emph{rate panel}, not an established pLGH-convergence panel.
\end{theorem}

\begin{proof}[Proof sketch]
The bridge $\Wfun$ sends the Krein-side propinquity bound $\Lprop$ to
the synthetic-side pLGH rate $d_{\mathrm{pLGH}}$. The functor is not
isometric ($\BigDeth$ and $\ell^{L}$ measure different physical
quantities), but the rate parameter $\gammajoint = O(\log\nmax/\nmax
+ T/\Nt)$ is the same object on both sides, hence the numerical panel
transports bit-identically.
\end{proof}

\begin{remark}[The substantive new content]\label{rem:t3_substantive}
All extant Mondino--S\"amann-style examples (arXiv:2504.10380~v4 \S~10)
are constructed geometrically from continuous spacetimes:\
Chruści\'el--Grant approximations of globally hyperbolic spacetimes,
sequences of warped products, etc. Theorem~\ref{thm:bit_exact_panel}
is the first quantitative pLGH-convergence rate panel in the MS
literature derived from \emph{operator-algebraic} input via the
bridge.
\end{remark}

\subsubsection{T4:\ Synthetic four-witness
readings}\label{sec:t4_synthetic_four_witness}

\begin{theorem}[T4:\ Synthetic four-witness
readings]\label{thm:synthetic_four_witness}
The Paper~42~\cite{paper42} four-witness Wick-rotation theorem
identifies six physical thermal-time witnesses (Bisognano--Wichmann
canonical, Hartle--Hawking at $M = 1, 2$, Sewell at $M = 1$, Unruh at
$a = 1, 2$) on the operator-algebraic side, all giving bit-identical
modular operators on the wedge KMS state. Via $\Wfun$, these six
physical witnesses give six geometric-time readings of the same
synthetic-Lorentzian wedge pre-length space, distinguished only by
the rapidity-rate normalization $\kappag$:

\begin{center}
\begin{tabular}{lcl}
\toprule
Witness & $\kappag$ & Synthetic time scale \\
\midrule
BW canonical & $1$ & $\ell^{L} = \tmod$ \\
Hartle--Hawking $M = 1$ & $1/4$ & $\ell^{L} = \tmod/4$ \\
Hartle--Hawking $M = 2$ & $1/8$ & $\ell^{L} = \tmod/8$ \\
Sewell $M = 1$ & $1/4$ & same as HH $M = 1$ \\
Unruh $a = 1$ & $1$ & same as BW canonical \\
Unruh $a = 2$ & $2$ & $\ell^{L} = 2\tmod$ \\
\bottomrule
\end{tabular}
\end{center}

Witnesses with the same $\kappag$ produce identical synthetic LPLS:\
the six physical instantiations on the operator side reduce to
\textbf{two synthetic-side equivalence classes}.
\end{theorem}

\begin{proof}[Proof sketch]
The bridge time separation $\ell^{L} = \kappag\cdot\tmod$ rescales
linearly with $\kappag$. The bit-exact algebra-action collapse of the
six witnesses on the Krein side (Paper~42~\cite{paper42} \S~8 Cor~8.1)
translates to a bit-exact metric-rescaling identification on the
synthetic side.
\end{proof}

\subsubsection{T5:\ Synthetic Pythagorean orthogonality with $1/\pi^{2}$
M1 signature}\label{sec:t5_synthetic_pythagorean}

\begin{theorem}[T5:\ Synthetic-side Pythagorean
orthogonality]\label{thm:synthetic_pythagorean}
Under the bridge $\Wfun$ applied to the GeoVac hemispheric wedge at
the Riemannian limit $\Nt = 1$, the BW-aligned thermal-time generator
$\Hloc = \Kalpha/\beta$ and the wedge-restricted Lorentzian Dirac
$\Dwedge$ correspond to two structurally orthogonal foliation
structures on the synthetic wedge pre-length space:\ a boost-orbit
foliation $\mathcal{F}_{\alpha}$ ($\Hloc$ generator, $\Pi_{W}$-even
per Paper~43~\cite{paper43} Lem~parity\_HS input (I1)) and a
spinor-bundle foliation $\mathcal{F}_{D}$ ($\Dwedge$ generator,
$\Pi_{W}$-odd per input (I2)).

The Paper~43~\cite{paper43} \S~10.2 Pythagorean closed-form residual
\[
\Frob{\Hloc - \Dwedge}^{2} =
\frac{\kappag^{2}\cdot S(\nmax)}{4\pi^{2}} + D(\nmax)
\]
translates to a synthetic-side decomposition of the wedge's geometric
data into a $\kappag^{2}$-scaling Hopf-base-measure contribution
($S(\nmax)/(4\pi^{2})$) and a $\kappag$-independent spinor-bundle
contribution ($D(\nmax)$). The $1/\pi^{2} = \mathrm{Vol}(\Sone)^{-2}$
master Mellin engine M1 signature
(Paper~32~\cite{paper32} \S~VIII \texttt{rem:pythagorean\_m1\_closure})
identifies as the inverse-squared measure of the bridge cover's
slab-period.
\end{theorem}

\begin{proof}[Proof sketch]
The algebraic identity is Paper~43~\cite{paper43} Thm~pythagorean\_formal.
The synthetic-side reading is by application of $\Wfun$:\ the
$\Pi_{W}$-grading on the Krein side translates to the orthogonal
foliation structure on the synthetic side via Gelfand duality, and the
$1/\pi^{2}$ prefactor on the operator-algebraic Pythagorean residual
IS the synthetic-side wedge's slab-period inverse-squared measure
($\beta = 2\pi$, $1/\beta^{2} = 1/(4\pi^{2})$).
\end{proof}

\subsubsection{T6:\ G2 metric-level closure for the GeoVac
wedge}\label{sec:t6_g2_wedge_closure}

\emph{[Descoped (see the Status note): the ``metric-equivalent
level'' claimed below quantifies over the degenerate Paper~45
seminorm; the metric-level closure is open pending the upstream
repair. The norm-resolvent ingredients from Paper~47 are
unaffected.]}

\begin{theorem}[T6:\ G2 metric-level \emph{route} for the GeoVac
hemispheric wedge (descoped --- conditional on the Paper~45 degenerate seminorm)]\label{thm:g2_wedge_closure}
For the GeoVac hemispheric wedge specifically, the de-compactification
limit $T \to \infty$ closes at the \emph{propinquity-equivalent
level on the synthetic side}:\ the synthetic wedge images
$\Wfun(\mathbb{X}^{K}_{\mathrm{wedge}}(k))$ at admissible-scaling
cutoffs converge in the Mondino--S\"amann pLGH topology to the
synthetic non-compact wedge object $\Wfun(\mathbb{X}^{K}_{\mathrm{wedge}}(\infty))$
(where $\mathbb{X}^{K}_{\mathrm{wedge}}(\infty)$ is the bridge image
of the Paper~47~\cite{paper47} non-compact carrier $\sthree\times\R_{t}$
wedge).

Equivalently, $d_{\mathrm{pLGH}}(\Wfun(\mathbb{X}^{K}_{\mathrm{wedge}}(k)),
\Wfun(\mathbb{X}^{K}_{\mathrm{wedge}}(\infty))) \to 0$ as $T(k) \to \infty$.

This is a \textbf{GeoVac-wedge-specific proposed route toward} $\mathrm{G2}$
(de-compactification limit, Paper~45~\cite{paper45} \S~1.4) at the
\textbf{metric-equivalent level on the synthetic side} --- \emph{descoped},
conditional on the descoped Krein seminorm (Paper~45); it is not an
established closure.
\end{theorem}

\begin{proof}[Proof sketch]
By Paper~47~\cite{paper47} \S~6, the three Krein-side carriers
(periodic compact, bounded Dirichlet, non-compact $\R_{t}$) all
converge in norm-resolvent sense to the same physical Lorentzian
Dirac as $T \to \infty$. Under the bridge $\Wfun$:\
$\Wfun(\mathbb{X}^{K}_{\mathrm{per}}(T))$,
$\Wfun(\mathbb{X}^{K}_{\mathrm{Dir}}(T))$, and
$\Wfun(\mathbb{X}^{K}_{\R})$ are well-defined Mondino--S\"amann
pre-length spaces (T1 applies). Norm-resolvent convergence on the
Krein side plus the bridge functoriality plus the (B2)+(B4) bridge
properties give pLGH convergence on the synthetic side.

The substantive content is that the synthetic-side bridge image
carries the structural content of metric-level convergence even
though the Krein-side propinquity $\Lprop$ is not defined on
non-compact carriers (the Paper~45 propinquity requires compact
temporal radius).
\end{proof}

\subsubsection{T6a:\ Synthetic three-carrier
identification}\label{sec:t6a_three_carrier_synthetic}

\begin{corollary}[T6a:\ Synthetic three-carrier
identification]\label{cor:three_carrier_synthetic}
The Paper~47~\cite{paper47} \S~6 three-carrier identification on the
operator-algebraic side translates via the bridge $\Wfun$ to a
synthetic-side three-carrier identification:\ the three synthetic
wedge images $(\Wfun(\mathbb{X}^{K}_{\mathrm{per}}),
\Wfun(\mathbb{X}^{K}_{\mathrm{Dir}}), \Wfun(\mathbb{X}^{K}_{\R}))$
are isometric as Mondino--S\"amann covered LPLS in the limit
$T \to \infty$.
\end{corollary}

\begin{proof}[Proof sketch]
The Paper~47 isometric embeddings $\widetilde{J}_{T}^{\mathrm{per}},
\widetilde{J}_{T}^{\mathrm{Dir}}: \Krein_{T}^{\bullet} \hookrightarrow
\Krein_{\R}$ commute with the M-diagonal topographies (which are
time-independent), hence descend via Gelfand duality to maps on the
synthetic side. The synthetic three-carrier identification is the
dual of the Krein-side identification.
\end{proof}

\subsection{Honest scope of the wedge application}\label{sec:wedge_scope}

The bridge construction itself
applies cleanly to every load-bearing GeoVac structure (Paper~42
four-witness, Paper~43 wedge + Pythagorean, Paper~45 panel, Paper~47
three-carrier). The headline content (T6 G2 metric-equivalent route
for the wedge) is a \emph{conditional} design --- it inherits the
descoped Paper~45 degeneracy, so the metric-level closure is
\emph{not} established (descoped), only a proposed route; the general
non-compact Latr\'emoli\`ere problem
remains open. The strong-form bridge to the enlarged substrate is
Q1' (Section~\ref{sec:open_q1prime}).

% =====================================================================
\section{Synthetic-Lorentzian compactness
theorems}\label{sec:compactness}
% =====================================================================

\subsection{Synthetic compactness of the GeoVac wedge
family}\label{sec:syn_compactness_wedge}

Theorem~\ref{thm:synthetic_compactness} (T2) provides the synthetic
compactness statement for the GeoVac wedge family. We restate it in
the context of synthetic-Lorentzian compactness theory.

\begin{theorem}[Synthetic compactness restated]\label{thm:syn_compactness_restated}
The class $\mathfrak{X}_{\mathrm{GeoVac\ wedge}}$ of bridge images of
GeoVac wedge Krein PPQMS at admissible-scaling cutoffs is
pre-compact in the Mondino--S\"amann pLGH topology of
Definition~\ref{def:plgh}.
\end{theorem}

This is a substantive new content for the Mondino--S\"amann literature:\
the wedge family is the first operator-algebraically derived class to
satisfy the Mondino--S\"amann pre-compactness criteria.

\subsection{The first quantitative pLGH panel from
operator-algebraic input}\label{sec:first_quant_plgh_panel}

By Theorem~\ref{thm:bit_exact_panel} (T3), the GeoVac wedge family
furnishes the first quantitative pLGH-convergence panel derived from
operator-algebraic input. The bit-exact numerical values
$\{2.0746, 1.6101, 1.3223\}$ across $(\nmax, \Nt) \in \{(2, 3),
(3, 5), (4, 7)\}$ provide a concrete benchmark for future
Mondino--S\"amann-style synthetic-Lorentzian numerical work.

\subsection{G2 metric-equivalent closure for the GeoVac
wedge (proposed/descoped)}\label{sec:g2_wedge_closure_unpacked}

The deepest \emph{proposed} content of the paper is
Theorem~\ref{thm:g2_wedge_closure} (T6) --- a conditional route, \emph{descoped}
(it rests on Paper~45's degenerate Krein-side $\Lprop$;\ see the Status
note). We unpack the structural significance:

\begin{itemize}
\item \textbf{The non-compact Latr\'emoli\`ere propinquity is not
in the published literature.} The original Latr\'emoli\`ere
propinquity~\cite{latremoliere2018, latremoliere_metric_st_2017} and
its 2025 hypertopology extension~\cite{latremoliere2025_hypertopology}
operate on compact (or proper, locally compact) quantum metric spaces.
A non-compact extension to the non-compact carrier $\sthree\times\R_{t}$
remains open. This is the published-open problem $\mathrm{G2}$
metric-level of Paper~45~\cite{paper45} \S~1.4.
\item \textbf{The bridge supplies a workaround at the wedge level.}
The Wick-rotation functor $\Wfun$ converts the operator-algebraic
non-compact propinquity question to a synthetic-Lorentzian pLGH
convergence question (Mondino--S\"amann pLGH on non-compact LPLS is
well-defined). On the repaired (non-degenerate) seminorm the
synthetic-side closure \emph{would} exist and the bridge \emph{would}
transport it back to the operator side --- but as it stands this rests
on Paper~45's degenerate $\Lprop$, so T6 is a \emph{proposed/conditional}
route, not an established closure (Status note).
\item \textbf{The closure is wedge-specific.} The bridge applies to
any Krein PPQMS, but the headline content (G2 closure) requires the
$2\pi$-periodicity of modular flow on the wedge and the bridge's
ability to extend along the admissible-scaling sequence to the
non-compact carrier. For more general (non-wedge) carriers, the
$2\pi$-periodicity does not hold and the closure does not extend.
\end{itemize}

\subsection{Synthetic three-carrier
identification}\label{sec:synthetic_three_carrier}

Corollary~\ref{cor:three_carrier_synthetic} (T6a) extends the
Paper~47~\cite{paper47} three-carrier identification from the
operator side to the synthetic side. This resolves an apparent
tension between Papers~43 and~45's different temporal-carrier choices
(bounded Dirichlet vs.\ periodic compact) at the synthetic level:\
both, plus the non-compact carrier, are isometric as covered LPLS in
the de-compactification limit.

\subsection{Why this is genuinely new content for Mondino--S\"amann
pLGH theory}\label{sec:why_new_for_ms}

The Mondino--S\"amann pLGH theory~\cite{mondino_samann2025_pointed}
develops the synthetic side independently of operator algebras. All
extant examples in arXiv:2504.10380~v4 \S~10 are constructed
geometrically:\ Chruści\'el--Grant approximations of globally
hyperbolic spacetimes, warped products, etc. The bridge construction
of this paper introduces a categorically new source of MS-style
examples:\ operator-algebraic truncations transported to the synthetic
side via the Wick-rotation functor. The bridge construction itself is
a genuinely new \emph{source} of MS-style examples; the downstream
headlines (T3 quantitative panel, T6 G2 route) are \emph{conditional}
on the bridge and descoped at the metric level (they inherit the
Paper~45 degeneracy), so they are proposed content, not established
new results, for the Mondino--S\"amann literature.

\subsection{Related synthetic-Lorentzian convergence
frameworks}\label{sec:related_syn_lor}

We position the bridge construction vs.\ adjacent synthetic-Lorentzian
convergence frameworks. The Sormani--Vega null
distance~\cite{sormani_vega2016} and Sakovich--Sormani timed
Gromov--Hausdorff~\cite{sakovich_sormani2024} operate via null-distance
metric-GH rather than via Mondino--S\"amann's
ε-net-of-causal-diamonds. Minguzzi--Suhr bounded Lorentzian metric
spaces~\cite{minguzzi_suhr2024} extend a complementary distinction
metric. Allen--Burtscher metric GH~\cite{allen_burtscher2022} and
M\"uller's Cauchy slabs~\cite{muller2022} provide alternative
frameworks. The bridge of this paper targets specifically the
Mondino--S\"amann pLGH framework on covered LPLS; a parallel bridge to
null-distance or timed-Hausdorff frameworks would be a separate
construction.

% =====================================================================
\section{Open questions}\label{sec:open_questions}
% =====================================================================

\subsection{Q1':\ Strong-form bridge on the enlarged
substrate}\label{sec:open_q1prime}

The most natural follow-on to the K$^{+}$-weak-form bridge of this
paper is the strong-form bridge extending $\Wfun$ to the Paper~46
\cite{paper46} Appendix~B enlarged substrate (chirality-flipping
generators $M^{\mathrm{flip}}$ with $\{\JL, M^{\mathrm{flip}}\} = 0$).

\begin{quote}
\textbf{Open question Q1' (strong-form Krein--MS bridge).} \emph{Does
the Wick-rotation functor $\Wfun$ extend to the enlarged substrate of
Paper~46~\cite{paper46} Appendix~B, and if so, does the off-orbit
super-additivity close via direct transport of the
Paper~42~\cite{paper42} four-witness theorem to MS time separation?}
\end{quote}

\subsubsection*{The substantive question that A.4'-A
surfaced}

On the K$^{+}$-weak-form bridge of this paper, off-orbit
super-additivity is structurally vacuous via
Lemma~\ref{lem:decomposition_o} (Case (iii) is empty under the
M-diagonal topography). On the strong-form bridge, the enlarged
topography makes Case (iii) non-empty:\ there are triples
$(\hat{\omega}_{x}, \hat{\omega}_{y}, \hat{\omega}_{z})$ such that
all three pairs are on different orbits with non-trivial
chirality-flipping content.

In this regime, the operator-algebraic super-additivity from
Paper~42's four-witness theorem becomes load-bearing:\ the on-orbit
equality $\ell^{L}(x, y) + \ell^{L}(y, z) = \ell^{L}(x, z)$
generalizes to a strict super-additivity $\ell^{L}(x, y) + \ell^{L}(y,
z) \le \ell^{L}(x, z)$ via the analog of the special-relativistic twin
paradox at the operator-algebraic level.

\subsubsection*{Q1' three-step decomposition}

\begin{itemize}
\item[\textbf{(Q1'.A)}] \textbf{Topography enlargement.} Extend the
topography $\Mcal^{L}$ to include $M^{\mathrm{flip}}$ generators;
verify the topography axioms (Definition~\ref{def:krein_topography})
hold for the enlarged structure. Substantive risk:\ the enlargement
may break Abelianness (chirality-flipping generators do not commute
with each other in general), in which case the Gelfand-spectrum
construction does not apply directly.
\item[\textbf{(Q1'.B)}] \textbf{Gelfand spectrum.} Verify
$\hat{\Mcal}^{L}_{\mathrm{enlarged}}$ is well-defined. If Abelianness
breaks:\ generalize to non-commutative Mondino--S\"amann pre-length
spaces (a substantial NCG-research target not currently in the
literature).
\item[\textbf{(Q1'.C)}] \textbf{Off-orbit super-additivity.}
Prove $\ell^{L}(x, y) + \ell^{L}(y, z) \le \ell^{L}(x, z)$ at
strict-inequality level via Paper~42 four-witness theorem transport
to off-axis modular flow.
\end{itemize}

\subsubsection*{Q1' effort estimate and recommendation}

Estimated effort:\ 3--6 months at sprint cadence. If Q1'.B forces the
non-commutative MS extension route, the timeline expands further
(the non-commutative MS pre-length space concept does not exist in
the published literature; constructing it would itself be a 6--12 month
NCG-mathematics target). Q1' is named here as the principal multi-month
follow-on; it is the natural target for a Sprint L3e-Q1' or for an
independent NCG-research program.

\subsection{G2 metric-level closure on the general (non-wedge)
carrier}\label{sec:g2_general}

Theorem~\ref{thm:g2_wedge_closure} (T6) proposes a descoped, conditional
route to G2 metric-level for the GeoVac wedge via the bridge image on the synthetic side.
The general non-compact Latr\'emoli\`ere propinquity problem
(Paper~45~\cite{paper45} \S~1.4 $\mathrm{G2}$-metric on arbitrary
non-wedge carriers) remains the multi-month NCG-mathematics frontier.
Paper~47~\cite{paper47} establishes the analog closure at the
norm-resolvent / spectral level on the general carrier; the
metric-level analog requires either:
\begin{itemize}
\item A non-compact Latr\'emoli\`ere propinquity extension (a
multi-month NCG-math target); or
\item A bridge to a non-wedge synthetic-Lorentzian framework that
admits metric-level pLGH convergence on the general carrier
(structurally distinct from the wedge bridge here).
\end{itemize}

\subsection{G3 cross-manifold $\Tcal_{\sthree} \otimes
\Tcal_{\mathrm{Hardy}}(S^{5})$}\label{sec:g3}

The cross-manifold extension of the Lorentzian Krein spectral triple
program is blocked at the NCG-framework level by the
Paper~24~\cite{paper24} \S~V four-layer Coulomb/HO asymmetry,
including the (layer iv) modular-Hamiltonian structure of the wedge
KMS state surfaced by Sprint~L2-F.1. This obstruction is orthogonal
to the present bridge:\ G3 is a structural extension target for the
NCG framework itself, not for the bridge to Mondino--S\"amann.

\subsection{G4 inner-factor calibration data}\label{sec:g4}

The W3 spectral-zeta candidate for inner-factor calibration data was
falsified (CLAUDE.md Sprint~W3, 2026-05-08:\ CKM Wolfenstein
candidate; three follow-up tracks ruled out PMNS, lepton masses, and
the mechanically-generated basis). No concrete proposals for
inner-factor calibration data remain in the literature. G4 is
orthogonal to the present bridge.

\subsection{Q2':\ Non-commutative Mondino--S\"amann pre-length
spaces}\label{sec:q2_prime}

If Q1'.B (Section~\ref{sec:open_q1prime}) forces breaking the
Abelianness of the enlarged topography, the bridge construction would
naturally generalize to non-commutative MS pre-length spaces. This is
a substantial NCG-research target not currently in the literature.

% =====================================================================
\section*{Acknowledgments}
% =====================================================================

This work is part of the GeoVac math.OA arc (Papers~29, 32, 38--47).
The bridge construction was opened by the Latr\'emoli\`ere
2512.03573~\cite{latremoliere2025_hypertopology} and Mondino--S\"amann
arXiv:2504.10380~v4~\cite{mondino_samann2025_pointed} December~2025
publications, with the Connes--Rovelli thermal-time
hypothesis~\cite{connes_rovelli1994} supplying the unifying
interpretation. Phase~A formalization was conducted via six sub-sprint
memos (Tier 3-Light diagnostic, A.2' Krein-lift, A.3' bridge
construction + concurrent-work re-check, A.4'-A G-B2 closure, A.4'-B
G-B4 closure, A.4'-C GeoVac wedge application) and synthesized at
Phase~A.5'. The paper draft is Phase~B of the Sprint~L3e-P3 program.

% =====================================================================
\begin{thebibliography}{99}
% =====================================================================

\bibitem[Allen \& Burtscher(2022)]{allen_burtscher2022}
B.~Allen, A.~Burtscher, ``Properties of the null distance and
spacetime convergence,'' \emph{Int. Math. Res. Not. IMRN} (2022),
no.~10, 7729--7808.

\bibitem[Bertozzini, Conti \& Lewkeeratiyutkul(2009)]{bertozzini_conti_lewkeeratiyutkul2009}
P.~Bertozzini, R.~Conti, W.~Lewkeeratiyutkul, ``Modular Theory,
Non-Commutative Geometry and Quantum Gravity,''
SIGMA 6 (2010), 067; arXiv:1007.4094.

\bibitem[Bisognano \& Wichmann(1976)]{bisognano_wichmann1976}
J.~J.~Bisognano, E.~H.~Wichmann, ``On the duality condition for
quantum fields,'' \emph{J. Math. Phys.} \textbf{17} (1976), 303--321.

\bibitem[Bizi, Brouder \& Besnard(2018)]{bizi_brouder_besnard2018}
N.~Bizi, C.~Brouder, F.~Besnard, ``Space and time dimensions of
algebras with applications to Lorentzian noncommutative geometry and
quantum electrodynamics,'' \emph{J. Math. Phys.} \textbf{59} (2018),
062303. arXiv:1611.07062.

\bibitem[Braun \& S\"amann(2025)]{braun_samann2025}
M.~Braun, C.~S\"amann, ``Gromov's reconstruction theorem and measured
Gromov--Hausdorff convergence in Lorentzian geometry,''
arXiv:2506.10852 (June~2025).

\bibitem[Camporesi \& Higuchi(1996)]{camporesi_higuchi1996}
R.~Camporesi, A.~Higuchi, ``On the eigenfunctions of the Dirac
operator on spheres and real hyperbolic spaces,''
\emph{J. Geom. Phys.} \textbf{20} (1996), 1--18.

\bibitem[Chamseddine \& Connes(1997)]{chamseddine_connes1997}
A.~H.~Chamseddine, A.~Connes, ``The spectral action principle,''
\emph{Comm. Math. Phys.} \textbf{186} (1997), 731--750.

\bibitem[Che, Perales \& Sormani(2025)]{che_perales_sormani2025}
M.~Che, R.~Perales, C.~Sormani,
``Gromov's compactness theorem for the intrinsic timed-Hausdorff distance,''
\emph{Differential Geometry and its Applications} \textbf{103} (2026), arXiv:2510.13069.

\bibitem[Connes(1994)]{connes1994}
A.~Connes, \emph{Noncommutative Geometry}, Academic Press, San Diego,
1994.

\bibitem[Connes \& Rovelli(1994)]{connes_rovelli1994}
A.~Connes, C.~Rovelli, ``Von Neumann algebra automorphisms and
time-thermodynamics relation in generally covariant quantum theories,''
\emph{Class. Quantum Grav.} \textbf{11} (1994), 2899--2918.

\bibitem[Connes \& van Suijlekom(2021)]{connes_vs2021}
A.~Connes, W.~D.~van~Suijlekom, ``Spectral truncations in
noncommutative geometry and operator systems,''
\emph{Comm. Math. Phys.} \textbf{383} (2021), 2021--2067.
arXiv:2004.14115.

\bibitem[Farsi \& Latr\'emoli\`ere(2024)]{farsi_latremoliere2024}
C.~Farsi, F.~Latr\'emoli\`ere, ``Collapse in noncommutative geometry
and spectral continuity,'' arXiv:2404.00240 (2024).

\bibitem[Farsi \& Latr\'emoli\`ere(2025)]{farsi_latremoliere2025}
C.~Farsi, F.~Latr\'emoli\`ere, ``Continuity for the spectral
propinquity of the Dirac operators associated with an analytic path
of Riemannian metrics,'' arXiv:2504.11715 (April 2025).

\bibitem[Franco \& Eckstein(2014)]{franco_eckstein2014}
N.~Franco, M.~Eckstein, ``An algebraic formulation of causality for
noncommutative geometry,''
\emph{Class. Quantum Grav.} \textbf{30} (2013), 135007 (and follow-ups).

\bibitem[Hartle \& Hawking(1976)]{hartle_hawking1976}
J.~B.~Hartle, S.~W.~Hawking, ``Path-integral derivation of black-hole
radiance,'' \emph{Phys.~Rev.~D} \textbf{13} (1976), 2188--2203.

\bibitem[Hekkelman \& McDonald(2024)]{hekkelman_mcdonald2024}
E.-M.~Hekkelman and E.~McDonald, ``A noncommutative integral on
spectrally truncated spectral triples, and a link with quantum
ergodicity,'' arXiv:2412.00628 (2024).

\bibitem[Ketterer(2026)]{ketterer2026}
C.~Ketterer, ``Convergence of Lorentzian spaces and curvature bounds
for generalized cones,'' arXiv:2605.11271 (May 2026).

\bibitem[Kubota(2026)]{kubota2026}
H.~Kubota, ``A Lorentzian coarea inequality,'' arXiv:2605.09101
(May 2026).

\bibitem[Kunzinger \& S\"amann(2018)]{kunzinger_samann2018}
M.~Kunzinger, C.~S\"amann, ``Lorentzian length spaces,''
\emph{Ann.~Global Anal.~Geom.} \textbf{54} (2018), 399--447.
arXiv:1711.08990.

\bibitem[Kostant(1969)]{kostant1965}
B.~Kostant, ``On the existence and irreducibility of certain series
of representations,'' \emph{Bull. Amer. Math. Soc.} \textbf{75}
(1969), 627--642.

\bibitem[Latr\'emoli\`ere(2015)]{latremoliere_metric_propinquity_2015}
F.~Latr\'emoli\`ere, ``The dual Gromov--Hausdorff propinquity,''
\emph{J. Math. Pures Appl.} \textbf{103} (2015), 303--351.

\bibitem[Latr\'emoli\`ere(2018)]{latremoliere2018}
F.~Latr\'emoli\`ere, ``The quantum Gromov--Hausdorff propinquity,''
\emph{Trans. Amer. Math. Soc.} \textbf{368} (2016), 365--411.
arXiv:1302.4058.

\bibitem[Latr\'emoli\`ere(2018/AF)]{latremoliere_pp_2018af}
F.~Latr\'emoli\`ere, ``Convergence of Cauchy sequences for the
covariant Gromov--Hausdorff propinquity,'' \emph{J. Math. Anal. Appl.}
\textbf{469} (2019), 378--404.

\bibitem[Latr\'emoli\`ere(2022)]{latremoliere_metric_st_2017}
F.~Latr\'emoli\`ere, ``The Gromov--Hausdorff propinquity for
metric spectral triples,'' \emph{Adv. Math.} \textbf{404} (2022),
Paper No.~108393. arXiv:1811.10843.

\bibitem[Latr\'emoli\`ere(2025)]{latremoliere2025_hypertopology}
F.~Latr\'emoli\`ere, ``The quantum Gromov--Hausdorff hypertopology
on the class of pointed proper quantum metric spaces,''
arXiv:2512.03573 (December~2025).

\bibitem[Leimbach(2024)]{leimbach_vs2024}
M.~Leimbach, ``Convergence of Peter--Weyl
truncations of compact quantum groups,'' to appear, \emph{J.~Noncommut.
Geom.} (2024); arXiv:2409.16698.

\bibitem[Marcolli \& van Suijlekom(2014)]{marcolli_vs2014}
M.~Marcolli, W.~D.~van~Suijlekom, ``Gauge networks in noncommutative
geometry,'' \emph{J. Geom. Phys.} \textbf{75} (2014), 71--91.
arXiv:1301.3480.

\bibitem[Minguzzi \& Suhr(2024)]{minguzzi_suhr2024}
E.~Minguzzi, S.~Suhr, ``Lorentzian metric spaces and their
Gromov--Hausdorff convergence,''
\emph{Lett.~Math.~Phys.} \textbf{114} (2024), Paper No.~73, arXiv:2209.14384.

\bibitem[Mondino, Ryborz \& S\"amann(2026)]{mondino_ryborz_samann2025}
A.~Mondino, V.~Ryborz, C.~S\"amann, ``Stability of synthetic timelike
Ricci bounds under $C^{0}$-limits and applications to impulsive
gravitational waves,'' arXiv:2605.03172 (May~2026).

\bibitem[Mondino \& S\"amann(2025)]{mondino_samann2025_pointed}
A.~Mondino, C.~S\"amann, ``Lorentzian Gromov--Hausdorff convergence
and pre-compactness,'' arXiv:2504.10380~v4 (December~2025).

\bibitem[M\"uller(2022)]{muller2022}
O.~M\"uller, ``Lorentzian Gromov--Hausdorff theory and finiteness
results,'' \emph{Gen.~Relativ.~Gravit.} \textbf{54} (2022), Art.~117;
arXiv:1912.00988.

\bibitem[Nieuviarts(2025)]{nieuviarts2025_v2}
G.~Nieuviarts, ``Emergence of Time from a Twisted Spectral Triple in
Almost-Commutative Geometry,'' arXiv:2512.15450 (December~2025, May~2026 v2).

\bibitem[Reed \& Simon(1978)]{reed_simon_iv}
M.~Reed, B.~Simon, \emph{Methods of Modern Mathematical Physics
Vol.~IV:\ Analysis of Operators}, Academic Press, New York, 1978.

\bibitem[Sakovich \& Sormani(2024)]{sakovich_sormani2024}
A.~Sakovich, C.~Sormani, ``Introducing various notions of distances
between space-times,'' (2024). arXiv:2410.16800.

\bibitem[Sewell(1982)]{sewell1982}
G.~L.~Sewell, ``Quantum fields on manifolds:\ PCT and gravitationally
induced thermal states,'' \emph{Ann.~Phys.~(NY)} \textbf{141} (1982),
201--224.

\bibitem[Sormani \& Vega(2016)]{sormani_vega2016}
C.~Sormani, C.~Vega, ``Null distance on a spacetime,''
\emph{Class. Quantum Grav.} \textbf{33} (2016), 085001. arXiv:1508.00531.

\bibitem[Strohmaier(2006)]{strohmaier2006}
A.~Strohmaier, ``On noncommutative and pseudo-Riemannian geometry,''
\emph{J. Geom. Phys.} \textbf{56} (2006), 175--195.

\bibitem[Unruh(1976)]{unruh1976}
W.~G.~Unruh, ``Notes on black-hole evaporation,'' \emph{Phys.~Rev.~D}
\textbf{14} (1976), 870--892.

\bibitem[van den Dungen(2016)]{vandungen2016}
K.~van den Dungen, ``Krein spectral triples and the fermionic action,''
\emph{Math. Phys. Anal. Geom.} \textbf{19} (2016), 4. arXiv:1505.01939.

% --- Internal GeoVac preprints ---

\bibitem[Loutey, internal Paper~24]{paper24}
J.~Loutey, ``The Bargmann--Segal Lattice: $\pi$-Free Discretization
of the Harmonic Oscillator on $S^5$,'' GeoVac internal preprint Paper~24
(2026), at
\texttt{papers/group3\_foundations/paper\_24\_bargmann\_segal.tex}.

\bibitem[Loutey, internal Paper~29]{paper29}
J.~Loutey, ``The GeoVac Hopf Graphs Are Ramanujan: Ihara Zeta,
Graph-RH, and a Scope Boundary on Selberg-on-Hydrogen,'' GeoVac internal preprint
Paper~29 (2026), at
\texttt{papers/group1\_operator\_algebras/paper\_29\_ramanujan\_hopf.tex}.

\bibitem[Loutey, internal Paper~32]{paper32}
J.~Loutey, ``The GeoVac Spectral Triple: A Synthesis Paper Locking the
Framework into the Marcolli--van~Suijlekom Gauge-Network Lineage,'' GeoVac internal preprint
Paper~32 (2026), at
\texttt{papers/group1\_operator\_algebras/paper\_32\_spectral\_triple.tex}.

\bibitem[Loutey, internal Paper~38]{paper38}
J.~Loutey, ``State-space Gromov--Hausdorff convergence of truncated
Camporesi--Higuchi spectral triples on $S^3$,'' GeoVac internal preprint Paper~38 (2026),
released to Zenodo 2026-05-07. At
\texttt{papers/group1\_operator\_algebras/paper\_38\_su2\_propinquity\_convergence.tex}.

\bibitem[Loutey, internal Paper~39]{paper39}
J.~Loutey, ``Tensor-product propinquity convergence on
$\Tcal_{S^{3}}^{\lambda_{a}}\otimes\Tcal_{S^{3}}^{\lambda_{b}}$,''
GeoVac internal preprint Paper~39 (2026), at
\texttt{papers/group1\_operator\_algebras/paper\_39\_tensor\_propinquity\_convergence.tex}.

\bibitem[Loutey, internal Paper~40]{paper40_unified}
J.~Loutey, ``State-space Gromov--Hausdorff convergence of spectral
truncations of compact Lie groups with bi-invariant metric,'' GeoVac
internal preprint Paper~40 (2026), at
\texttt{papers/group1\_operator\_algebras/paper\_40\_unified\_propinquity\_convergence.tex}.

\bibitem[Loutey, internal Paper~42]{paper42}
J.~Loutey, ``Tomita--Takesaki modular structure on truncated $SU(2)$
spectral triples: four-witness Wick-rotation literal identification
at finite cutoff,'' GeoVac internal preprint Paper~42 (2026), at
\texttt{papers/group1\_operator\_algebras/paper\_42\_modular\_hamiltonian\_four\_witness.tex}.

\bibitem[Loutey, internal Paper~43]{paper43}
J.~Loutey, ``Lorentzian extension of the four-witness Wick-rotation
theorem on truncated $S^3 \times \mathbb{R}$ spectral triples at finite
cutoff,'' GeoVac internal preprint
Paper~43 (2026), at
\texttt{papers/group1\_operator\_algebras/paper\_43\_lorentzian\_extension.tex}.

\bibitem[Loutey, internal Paper~44]{paper44}
J.~Loutey, ``Operator-system extension of the BBB Krein spectral
triple at finite cutoff: propagation number and envelope dependence
on the Camporesi--Higuchi spinor bundle,'' GeoVac internal preprint Paper~44 (2026),
at
\texttt{papers/group1\_operator\_algebras/paper\_44\_lorentzian\_operator\_system.tex}.

\bibitem[Loutey, internal Paper~45]{paper45}
J.~Loutey, ``Lorentzian propinquity on truncated SU(2)
$\times$ U(1)$_{T}$ Krein spectral triples: a degeneracy theorem,'' GeoVac internal
preprint Paper~45 (2026), at
\texttt{papers/group1\_operator\_algebras/paper\_45\_lorentzian\_propinquity.tex}.

\bibitem[Loutey, internal Paper~46]{paper46}
J.~Loutey, ``Strong-form Lorentzian propinquity on
truncated SU(2) $\times$ U(1)$_{T}$ Krein spectral triples,''
GeoVac internal preprint Paper~46 (2026), at
\texttt{papers/group1\_operator\_algebras/paper\_46\_strong\_form\_lorentzian\_propinquity.tex}.

\bibitem[Loutey, internal Paper~47]{paper47}
J.~Loutey, ``Norm-resolvent convergence to the non-compact Lorentzian
Krein spectral triple on $\sthree\times\R_{t}$ and the two-rate
hybrid limit,'' GeoVac internal preprint Paper~47 (2026), at
\texttt{papers/group1\_operator\_algebras/paper\_47\_two\_rate\_hybrid\_convergence.tex}.

\end{thebibliography}

\end{document}
